%% ugp_theory_overview.tex — UGP Physics: A Narrative Overview
%% This file is included in the P00 survey guide as an introductory chapter.
%% Maintained separately for ease of updates.
%% Last updated: 2026-06-10
%%
%% DRAFT — for author review. Not final.

% ====================================================================
\clearpage
\section{The UGP Framework: Foundations, Structure, and Results}
\label{sec:overview}
% ====================================================================

\subsection{Introduction}
\label{subsec:overview-intro}

The Standard Model of particle physics has approximately 25 free parameters ---
coupling constants, fermion masses, mixing angles, the cosmological constant ---
whose numerical values have no theoretical origin within the Standard Model itself.
The \textbf{UGP Physics programme}, the quantitative branch of the Reflexive Reality programme~\cite{SpivackResearchProgrammes}, addresses this
directly: starting from one foundational principle --- \textbf{PSC (Perfect
Self-Containment)} --- with MDL (Minimum Description Length) as its operational
expression (Presentation Invariance, the Layer~II PSC axiom; PI $\equiv$ MDL,
machine-certified), and locality and symmetry as structural requirements, it
uniquely selects a continuum field theory whose internal arithmetic
reproduces the SM particle spectrum, all three gauge coupling constants, all fermion
masses, the electroweak mixing angle, CKM mixing parameters, and now gravitational
and quantum-mechanical structure --- without a single free parameter.  A discrete
algebraic certificate ($f_{\rm MDL}$, the CMCA) proves and compresses that
arithmetic; it is not the physical substrate.

What distinguishes UGP from prior unification programmes is that it works in the
opposite direction: rather than postulating a larger symmetry group that contains
the SM as a subgroup, it derives the SM gauge structure \emph{from below}, from the
symmetry and orbit structure of a continuum field whose unique selection is forced
by information-theoretic compressibility.  The physical substrate is not invented
to fit the data; it is the unique solution to a self-consistency equation.
The central object is the \textbf{$\Phi_{\rm MDL}$ field}, a
$\mathbb{Z}_7$-symmetric Klein--Gordon gauge field whose internal discrete symmetry
group is $F_{21} = \mathbb{Z}_7 \rtimes \mathbb{Z}_3$ --- the unique non-abelian
group of order~21, a subgroup of $\mathrm{SU}(3)$ --- forced by minimum description
length (MDL) minimality.  Stable topological kink defects of $\Phi_{\rm MDL}$ are
the elementary particles of the Standard Model.  A discrete algebraic certificate
--- the MDL-minimal $\mathbb{Z}_7$ cellular automaton $f_{\rm MDL}$ --- certifies
the internal symmetry structure of $\Phi_{\rm MDL}$ and generates all SM quantum
numbers from its orbit arithmetic; it is not the physical substrate.  Paper~P34
formalises the abstract substrate architecture as the Möbius triple $(A,e,[D])$,
with D1--D5 constraints governing transputation and the Born rule.  Quantum
mechanics, special relativity, and gravity all emerge from the field's structure;
none is postulated.

Results throughout the programme are certified at multiple epistemic levels.  Many
are \emph{machine-certified} in Lean~4 with zero \texttt{sorry} and the standard
Mathlib axiom signature --- a level we denote $\mathbf{A}_{\rm Lean}$.  Others are
analytically derived ($\mathbf{A/D}$) or numerically confirmed against PDG
measurements.  This overview traces the logical arc from substrate selection to the
complete unified framework, citing the relevant papers for technical detail.

% ====================================================================
\subsection{The Substrate: MDL Minimality and the Unique Selection of \texorpdfstring{$\Phi_{\rm MDL}$}{Phi\_MDL}}
\label{subsec:overview-substrate}

The framework begins with a question: among all field theories, which one minimizes
the total description length of the laws governing a self-contained universe?
The \textbf{Minimum Description Length (MDL)} principle --- the operational
expression of PSC's Presentation Invariance (PI $\equiv$ MDL, Layer~II PSC axiom,
machine-certified)~\cite{SpivackSM} --- formalises self-containment as a
Kolmogorov-complexity criterion.
Combined with locality and symmetry as structural requirements, MDL minimality
uniquely selects the $\mathbb{Z}_7$-symmetric sine-Gordon field $\Phi_{\rm MDL}$:
among all $\mathbb{Z}_N$-symmetric fields with BPS kinks and exact Lorentz invariance,
$\Phi_{\rm MDL}$ has the smallest Kolmogorov complexity~\cite{SpivackCMCA,SpivackPhiMDLField}.
(MDL is not an independent foundational axiom: PI is the Layer~II PSC axiom that
selects $N_{\rm gen}=3$, and MDL is its physics-vocabulary expression.  `Why MDL?'
reduces to `Why PSC?' --- the foundational question answered by the NEMS programme.)
The NEMS physics synthesis situating GTE within the broader NEMS programme is in \cite{SpivackPhysicsPortal}.

\paragraph{The physical substrate.}
The \textbf{$\Phi_{\rm MDL}$ field} is the physical substrate of the universe.
It is a real-valued $\mathbb{Z}_7$-symmetric Klein--Gordon gauge field evolving on
$\mathbb{R}^{3+1}$:
\[
  \mathcal{L} = \tfrac{1}{2}(\partial_\mu\Phi)^2 - V(\Phi),
  \qquad
  V(\Phi) = \frac{m_\varphi^2}{49}(1-\cos 7\Phi),
\]
with field mass $m_\varphi = m_\tau = 1776.86$~MeV, derived from the $F_{21}$
leptonic-sector self-consistency condition (SCC)~\cite{SpivackCompleteness}.
Its seven global minima $\Phi_k = 2\pi k/7$, $k \in \mathbb{Z}_7$, define the
vacuum manifold; the topological winding number $Q = (7/2\pi)\int\partial_x\Phi\,dx$
takes values in $\mathbb{Z}_7$ and furnishes a superselection charge.
Stable topological kink defects winding between adjacent vacua are identified with
the elementary particles of the Standard Model.  Exact Lorentz invariance
($\omega^2 = k^2 + m^2$) is machine-certified in Lean~4
($\mathbf{A}_{\rm Lean}$, \lean{poincare\_invariance\_of\_kg})~\cite{SpivackPhiMDLField}.

\paragraph{Scale hierarchy: substrate vs.\ particle.}
\begin{center}
\small
\begin{tabular}{@{}p{4.5cm}p{5.5cm}p{4.0cm}@{}}
\toprule
\textbf{Object} & \textbf{Length scale} & \textbf{Energy scale} \\
\midrule
CMCA fine-end reading (working hypothesis $a \approx \ell_{\rm Pl}$) & ${\sim}10^{-33}$~cm ($\ell_{\rm Pl}$) & ${\sim}10^{19}$~GeV \\
CMCA matter-sector reading (MDL-saturated; derived, conditional --- see below) & $a = \hbar c/\Lambda_{\rm GTE} \approx 0.097$~fm $\approx 10^{-14}$~cm & $\Lambda_{\rm GTE} \approx 2.0$~GeV \\
$\Phi_{\rm MDL}$ kink width & ${\sim}1/m_\varphi \approx 10^{-13}$~cm & $m_{\rm kink} = 290$~MeV \\
SM particles (Compton scale) & $10^{-13}$--$10^{-16}$~cm & MeV--TeV \\
\bottomrule
\end{tabular}
\end{center}
The CMCA, as a discrete algebraic certificate, carries no intrinsic lattice spacing: its algebraic content holds at every resolution (Algebraic Descent), and no finite-resolution CA is the physical substrate.  This is the physical content of the no-CA-replica theorem: a finite-resolution CA at any $M$ cannot exactly reproduce the smooth Compton-scale kink profile because the kink extends over $\sim M^{-1}$ lattice sites in the continuum limit.  A lattice spacing therefore belongs to a \emph{reading} --- a calibration of the certificate against $\Phi_{\rm MDL}$ observables --- and two readings are in use.  Under the conventional \emph{fine-end working hypothesis} $a \approx \ell_{\rm Pl}$ (a working hypothesis, not a GTE derivation), the BPS kink width $\sim 1/m_\varphi \sim 10^{-13}$~cm is twenty orders of magnitude above the lattice scale, spanning $M_{\rm Pl}/M_{\rm kink} = 21^{10}\cdot 7^{9}/2^{4} \approx 4.2\times10^{19}$ cells: kinks are Compton-wavelength-scale collective excitations of a much finer lattice, precisely as phonons are macroscopic excitations of an atomic lattice.  Under the MDL-saturated \emph{matter-sector reading} $a = \hbar c/\Lambda_{\rm GTE} \approx 0.097$~fm --- derived rather than postulated (Tape Saturation Theorem, $\CatAD$, conditional on the named Compton-Support Criterion, which is proven minimal; P50) --- the same kink spans exactly $\lambda_C/a = |\mathbb{Z}_7| = 7$ cells, and the substrate-regulated lattice campaigns of P42 operate at this reading.  The separation of the two readings is itself GTE arithmetic: the Planck-hierarchy monomial $M_{\rm Pl}/\Lambda_{\rm GTE} = 3^{10}\cdot 7^{18}/2^{4} \approx 6.0\times10^{18}$.  The Algebraic Lifting Theorem bridges the certificate to the Compton scale without requiring a simulation of the Compton scale from first principles.

\paragraph{The discrete algebraic certificate.}
At the algebraic level, the orbit structure of $\Phi_{\rm MDL}$ is encoded by a
compact discrete object: the \textbf{MDL-minimal $\mathbb{Z}_7$ update rule}
$f_{\rm MDL}$, defined over $\mathrm{GF}(7)$ as
$f_{\rm MDL}(L,C,R) = C + R - CR - LCR$, whose binary restriction is Rule~110.
$f_{\rm MDL}$ is \emph{not} the physical substrate; it is the
\textbf{algebraic certificate} --- a compact discrete representation of the
symmetry structure, orbit arithmetic, and generation sequence of $\Phi_{\rm MDL}$.
The $f_{\rm MDL}$ orbit over $\mathbb{Z}_7$ generates all SM quantum numbers via the
GTE cascade (\S\ref{subsec:overview-gte}).

The complete discrete algebraic certificate is the \textbf{Three-Layer Chiral
Minkowski CA (CMCA)}~\cite{SpivackCMCA}: Rule~110 (right chiral layer) plus
Rule~124 (left chiral layer) plus an inner $\tau_c$ Rule-110 clock per cell.
The CMCA certifies all five structural features of $\Phi_{\rm MDL}$: $\mathbb{Z}_7$
generation orbit ($N_{\rm gen}=3$), V$-$A chirality, Turing universality, and both
observer-level and dynamics-level special relativity.  Its MDL complexity
$K_{\rm CA} = 19$ bits saturates the construction-class lower bound
($\mathbf{A}_{\rm Lean}$: \lean{CMCAMDLMinimality.lean})~\cite{SpivackCMCA}.

\paragraph{The Three-Tape CMCA and 3+1D emergence.}
The single-tape CMCA is a 1+1D algebraic certificate.  The \textbf{Three-Tape CMCA}
(P45~\cite{SpivackThreeTapeCMCA}) extends this to three independent 1+1D CMCAs
sharing a single outer clock period $\tau_c^{\rm out}$.  The
\textbf{Dimensional Protocol Principle (DPP)}, machine-certified in Lean~4
($\mathbf{A}_{\rm Lean}$), states that this shared clock is \emph{both necessary
and sufficient} for 3+1D dynamics: without it the three tapes evolve as independent
1+1D systems; with it they exhibit a single 3+1D Minkowski spacetime.  The DPP
provides the first machine-verified derivation of why spacetime is $3+1$-dimensional
within the GTE framework.  Temporal synchronisation is the causal prerequisite for
spatial geometry --- the shared clock \emph{creates} space.

The three-tape architecture additionally produces all 33 SM charged-current vertices
via $\mathbb{Z}_7$ winding conservation, derives baryon number as a topological
charge, and recovers colour confinement from the MDL description-length gap
$\Delta K = \log_2 9 \approx 3.17$ bits forbidding free quarks (all $\mathbf{A}_{\rm Lean}$,
zero sorry, P45/P46~\cite{SpivackGTEPolynomialUFT}).

The three-tape gravity mechanism has been closed at $\mathbf{A/D}$ (CatAD).  The Physical Minimum Description Length (PMDL)
Poisson equation $\nabla^2\Phi = G_{\rm eff}\,p(w_x,w_y,w_z)$ is $\mathbf{A}_{\rm Lean}$;
the continuum-limit force law is derived analytically via the 3D Poisson Green's function
and multipole expansion.  For a source of Algebraic-Lifting radius $\sigma_{\rm AL}$,
the exact force is
\[
  F(b) \;=\; \frac{G_{\rm eff} M}{4\pi b^2}\,
  \Bigl[1 - \tfrac{\sigma_{\rm AL}^2}{2b^2} + O\!\bigl(\sigma_{\rm AL}/b\bigr)^4\Bigr]
  \;\xrightarrow{b/\sigma_{\rm AL}\to\infty}\;
  \frac{G_{\rm eff} M}{4\pi b^2},
\]
recovering exactly Newtonian $F\propto r^{-2}$.  Prior discrete-simulation measurements
($F\propto b^{-2.19\ldots-2.30}$) are finite-source-size artifacts: the exponent
converges monotonically to $-2.000 \pm 0.005$ as $b/\sigma_{\rm AL}\to\infty$, confirmed
numerically.  Three equivalent formulations are established: gradient-kick (Level~1
computational), metric-perturbation/GR (Level~2 geometric), and equivalence-principle
(deepest mechanical principle), with the Level~2 metric formulation isomorphic to the
Level~1 result in the Newtonian limit.

\paragraph{The two-level architecture.}
The GTE/$\Phi_{\rm MDL}$ framework operates at two
levels~\cite{SpivackCompleteness,SpivackCMCA}:
\begin{description}[leftmargin=2.5cm,labelwidth=2.5cm]
\item[\textit{Level~1}:] The CMCA --- a 1+1D discrete algebraic certificate.  Carries
  all five Level-1 structural features.  This is \emph{not} the physical universe; it
  is a mathematical proof system for the structure of $\Phi_{\rm MDL}$.
\item[\textit{Level~2}:] $\Phi_{\rm MDL}$ in $3+1$ dimensions --- the physical
  continuum field.  Quantum mechanics, Standard Model structure, and General
  Relativity all emerge at this level.  The continuum limit of the CMCA is
  identified with $\Phi_{\rm MDL}$ ($\mathbf{A}_{\rm Lean}$:
  \lean{CMCAContinuumLimit.lean}).
\end{description}
Orthogonal to this resolution-graded level scheme, the framework's objects occupy
a three-rung \emph{ontological ladder} naming their roles --- the
\textbf{mathematical substrate} (the UGP/GTE arithmetic: what selects), the
\textbf{certificate} (the 19-bit polynomial, $f_{\rm MDL}$, and the CMCA tape:
what certifies), and the \textbf{physical substrate} ($\Phi_{\rm MDL}$: what
exists) --- with the normative statement given in
P48~\cite{SpivackGTECompleteFramework}.

\paragraph{The GTE-Möbius architecture specification (P34).}
Paper~P34~\cite{SpivackGTEMobius} formalises the \emph{abstract architecture} shared
by both levels as the Möbius triple $(A,\,e,\,[D])$:
\begin{itemize}[noitemsep,topsep=2pt]
  \item \textbf{$A$} --- the GTE arithmetic carrier ($\mathbb{Z}_7$ winding dynamics,
    P28).  At Level~2, $A$ is realised as the PSC-admissible \emph{kink sector} of
    $\Phi_{\rm MDL}$ (topological charges $Q \in \{0,1,2,3\}$).
  \item \textbf{$e$} --- the self-encoding map conferring Turing completeness
    (P30/P40; GF(7) polynomial route and Cook's Rule~110 construction).
  \item \textbf{$[D]$} --- the class of PSC-consistent coherence measures satisfying
    five constraints D1--D5 (non-negativity, PSC invariance, non-computability,
    selection completeness, Born consistency).  $[D]$ realises \emph{transputation}:
    non-computable, observer-free outcome selection at the computation/transputation
    boundary.
\end{itemize}
The Möbius name reflects a single-surface architecture: computation (deterministic
$f_{\rm MDL}$ evolution) and transputation ($[D]$-selection) share one substrate
without a sharp syntactic boundary.  P34 is the \textbf{axiomatic foundation
paper} for Group~7 (P34--P46): it specifies \emph{what kind of object} underlies the
SM derivations; P35 synthesises the bidirectional correspondence; P42--P43 complete
the $\Phi_{\rm MDL}$ field theory with D1--D5 as explicit axioms.

\begin{tcolorbox}[colback=green!3,colframe=green!50!black,
  title={\textbf{How the substrates and CA models relate}},
  fonttitle=\bfseries,left=6pt,right=6pt,top=5pt,bottom=5pt]
\small
\begin{center}
\begin{tabular}{@{}p{0.22\textwidth}p{0.30\textwidth}p{0.38\textwidth}@{}}
\toprule
\textbf{Object} & \textbf{Role} & \textbf{Key papers} \\
\midrule
$\Phi_{\rm MDL}$ field &
  \textbf{Physical substrate} (Level~2): the universe &
  P42, P43, P38, P44 \\
CMCA (3-layer) &
  Level-1 \emph{algebraic certificate}: chiral null structure, SR clock &
  P41 \\
Three-Tape CMCA &
  Level-1 certificate extended to 3+1D via shared clock (DPP) &
  P45 \\
$f_{\rm MDL}$ / Rule~110 &
  Minimal $\mathbb{Z}_7$ polynomial shadow of $A$; orbit arithmetic &
  P28, P40, P46 \\
$(A,e,[D])$ triple &
  Abstract specification; D1--D5 axioms for QM/transputation &
  P34 \\
\bottomrule
\end{tabular}
\end{center}
\noindent\textbf{Final substrate:} $\Phi_{\rm MDL}$ in $3{+}1$D.  \textbf{Not the
substrate:} any finite CA, including $f_{\rm MDL}$, Rule~110, or the CMCA
(no-CA-replica theorem, P43).  \textbf{Relation:} discrete objects \emph{certify}
and \emph{compress} the arithmetic of $\Phi_{\rm MDL}$; lifting/descent theorems
(P34, P35) connect the levels.
\end{tcolorbox}

Results flow upward from discrete to physical via the \textbf{Algebraic Lifting
Theorem} (\lean{algebraic\_lifting\_theorem}, $\mathbf{A}_{\rm Lean}$) and downward
from continuum to discrete via the \textbf{Algebraic Descent Theorem}
(\lean{AlgebraicDescentTheorem.lean}, $\mathbf{A}_{\rm Lean}$): every algebraic and
topological property of $\Phi_{\rm MDL}$ depending only on the $F_{21}$ internal
symmetry --- colour confinement, three generations, asymptotic freedom, strong~CP
--- holds for $f_{\rm MDL}$ at every lattice resolution $M \geq 1$.

\paragraph{The Single-Source Principle.}
The GTE polynomial $p(L,C,R) = C + R - CR - LCR$ over $\mathrm{GF}(7)$, which
requires exactly 19~bits to specify, simultaneously plays \emph{five} physical
roles with zero additional specification cost per role
(P46~\cite{SpivackGTEPolynomialUFT}; all roles $\mathbf{A}_{\rm Lean}$):
(1)~\emph{spatial cellular-automaton dynamics} --- Rule~110 Turing-universal
Class-4 behaviour;
(2)~\emph{$\mathbb{Z}_7$ gauge vertex conservation} --- Standard Model
charged-current vertices;
(3)~\emph{gravitational coupling} --- the PMDL Poisson source term
$\nabla^2\Phi = G_{\rm eff}\,p(w_x,w_y,w_z)$;
(4)~\emph{quantum entanglement generator} (negativity $\mathcal{N} = 0.38$, Bell nonlocality $S = 2.4459$, $87\%$ of Tsirelson bound);
(5)~\emph{baryon number}: $B = (1/3)\sum \chi_q(w_j)$ topological charge, $B = 1/3$ per quark from $N_{\rm tapes} = 3$
($\mathbf{A}_{\rm Lean}$: \lean{gte\_baryon\_number\_topological\_charge}).
This constitutes the \textbf{Single-Source Principle}: all of spacetime, matter,
gravity, and quantum mechanics emerge from a single 19-bit description.
The MDL uniqueness of this polynomial is proved in Lean~4: it is the unique
degree-3 multilinear object over $\mathrm{GF}(7)$ satisfying $\mathbb{Z}_7$
symmetry, Rule~110 restriction, and 19-bit MDL minimality.

\paragraph{MDL uniqueness of $\mathbb{Z}_7$.}
The integer $N = 7$ is uniquely forced by three independent arguments:
(a)~MDL-minimality: $b_0 = 7N_{\rm gen}/3 = 7$ is a positive integer only for
$N=7$, by the Algebraic Necessity Theorem
($\mathbf{A}_{\rm Lean}$: \lean{algebraic\_necessity\_master\_bundle});
(b)~the Frobenius prime identity $|\mathbb{Z}_7| = |\mathbb{Z}_3|^2 - |\mathbb{Z}_3|
+ 1 = 7$ ($\mathbf{A}_{\rm Lean}$: \lean{FrobeniusPrimeIdentity.lean});
(c)~the gravitational hierarchy formula
$M_{\rm Pl}/m_\tau = |F_{21}|^{10}|\mathbb{Z}_7|^7/2$ at $0.040\%$ precision
requires $|\mathbb{Z}_7| = 7$~\cite{SpivackCompleteness}.
Any other value of $N$ fails at least one of these three constraints.

The five conditions that certify the orbit structure as uniquely admissible are
formalised in the Perfect Self-Containment (PSC) framework: renormalisability (RC),
structural stability (NM*), anomaly cancellation (AC), topological viability (TV),
and spectral asymmetry (SA).  Among $34{,}560$ candidate dynamical systems, exactly
12 pass Layer~I hard PSC filters ($0.03\%$); every Layer~I survivor carries
$N_{\rm gen}=3$ ($\CatAL$, zero sorry:
\lean{psc\_enumeration\_forces\_ngen\_3}).  Layer~II Presentation Invariance selects
the unique MDL-minimum survivor at ridge level $n{=}10$, the GTE Lepton Seed
$(1,\,73,\,823)$, proved in
Lean~4~\cite{SpivackPSCConcordance,SpivackCompUniversality,SpivackGTEUnification}.
(The Layer~II claim that Presentation Invariance selects the minimum-$K$ survivor
among the 12 PSC-admissible gauge theories remains to be Lean-certified.)
The Unified Rigidity Theorem~\cite{SpivackUnifiedRigidity} (P12) provides the
synthesis proof that these three independent selection mechanisms --- PSC Layer~I/II
theory-space selection, UGP arithmetic seed uniqueness, and the dynamical basin
structure --- all converge on the same canonical branch, with every assumption
explicitly typed and audited.

The \textbf{Law~$=$~Description~$=$~Execution identity} holds for the discrete
certificate: $f_{\rm MDL}$ is simultaneously (a)~the GTE evolution law,
(b)~a compressed description encoding the 14 nonzero SM neighbourhoods, and
(c)~the executing CA whose binary sublayer is Rule~110.  Machine-certified:
\lean{fmdl\_law\_description\_execution} in
\lean{FMDLClassification.lean} ($\mathbf{A}_{\rm Lean}$, zero sorry)~\cite{SpivackReflexiveLaw}.

\begin{tcolorbox}[colback=blue!3,colframe=blue!50!black,
  title={\textbf{The role of cellular automata in the GTE framework}},
  fonttitle=\bfseries,left=6pt,right=6pt,top=5pt,bottom=5pt]
\small
The GTE framework involves cellular automata in a specific and limited role
that differs from their role in other CA-based physics proposals.

\textbf{What $f_{\rm MDL}$ / Rule~110 / the CMCA are:}
\begin{itemize}[noitemsep,topsep=2pt]
  \item The \emph{algebraic certificate}: a compact discrete representation of
    $\Phi_{\rm MDL}$'s internal symmetry structure and orbit arithmetic.
  \item The source of all SM generation quantum numbers via its $\mathbb{Z}_7$ orbit.
  \item A structure whose binary restriction (Rule~110) is forced by the SM
    generation sequence --- proved in Lean~4 with zero sorry.
  \item The \emph{mathematical proof system} for why $\Phi_{\rm MDL}$ is the correct
    continuum theory, not an alternative to it.
\end{itemize}
\textbf{What they are not:}
\begin{itemize}[noitemsep,topsep=2pt]
  \item NOT the physical substrate of the universe ($\Phi_{\rm MDL}$ is).
  \item NOT the source of quantum mechanics (which emerges from $\Phi_{\rm MDL}$
    at the continuum Level~2).
  \item NOT a replacement for $\Phi_{\rm MDL}$: P43 proves
    (\lean{no\_finite\_ca\_exact\_lorentz\_replica}, $\mathbf{A}_{\rm Lean}$) that
    no finite-resolution CA, quantum CA, or rewriting system can exactly replicate
    $\Phi_{\rm MDL}$'s Lorentz invariance.  The Nyquist residual
    $\varepsilon_0(M) = \pi^2/(3M^2) > 0$ for all finite $M$; only the continuum
    limit $M \to \infty$ recovers exact Lorentz symmetry, uniquely identifying
    $\Phi_{\rm MDL}$~\cite{SpivackCompleteness}.
\end{itemize}
\end{tcolorbox}

\paragraph{The computational/algebraic reframe.}
The GTE framework realises a scientific reframe that is distinct from all prior
computational-universe proposals.  The physical substrate is the continuum field
$\Phi_{\rm MDL}$; no finite CA is the substrate (no-CA-replica theorem,
\lean{no\_finite\_ca\_exact\_lorentz\_replica}, $\mathbf{A}_{\rm Lean}$).  What is
fundamental is the \emph{algebraic structure}: $\mathbb{Z}_7$, $F_{21} =
\mathbb{Z}_7\rtimes\mathbb{Z}_3$, and the GTE orbit arithmetic.  Three consequences
follow at the foundational level.  First, the Law\,$=$\,Description\,$=$\,Execution
identity (\lean{fmdl\_law\_description\_execution}, $\mathbf{A}_{\rm Lean}$,
zero sorry) certifies that $f_{\rm MDL}$ is simultaneously the evolution law, a
compressed description, and an executing CA --- the law of physics and its
description are the same object.  Second, the algebraic structure certifying
$\Phi_{\rm MDL}$ is itself computationally universal: Rule~110's Turing universality
(Cook's theorem, 2004) transfers to $\Phi_{\rm MDL}$ via the proved stepwise kink
simulation, and the $\mathrm{GF}(7)$ polynomial representation additionally shows, via
its center-cell NAND identity, that the kink update rule realizes any finite Boolean
function algebraically --- a finite functional-completeness result, distinct from and
not sufficient for Turing completeness on its own.  Third, the geodesic computation
principle (\S\ref{subsec:overview-gravity})
identifies physical trajectories with minimum-description-length paths: gravity,
inertia, and quantum measurement are all manifestations of MDL minimality in
the algebraic structure.  These three results together constitute a reframe in which
computation is the \emph{deep structure}, not the medium, of physical reality.

% ====================================================================
\subsection{The Generative Triple Evolution (GTE) Framework}
\label{subsec:overview-gte}

The orbit arithmetic of $f_{\rm MDL}$ organises into canonical integer triples
$(a,\,b,\,c;\,g)$ --- the \textbf{Generative Triples} --- one per Standard Model
particle species.  Three structural constants emerge from the orbit:
\[
  N_{\rm gen} = 3, \qquad N_{\rm fam} = 5, \qquad c_H = 13.
\]

\emph{All three are derived, not assumed.}  $N_{\rm gen} = 3$ follows from the
Garden-of-Eden orbit depth: the first-generation beable has no $f_{\rm MDL}$
predecessor, forcing the orbit to close after exactly three generations
(Lean: \lean{fmdl\_gen1\_is\_garden\_of\_eden}, $\mathbf{A}_{\rm Lean}$).
$N_{\rm fam} = 2^{N_{\rm gen}} - N_{\rm gen} = 5$ follows from the $\mathbb{Z}_5$
ring structure of the generation orbit (Lean:
\lean{gte\_master\_formula\_complete}, $\mathbf{A}_{\rm Lean}$).
$c_H = 2^{N_{\rm gen}+1} - N_{\rm gen} = 13$ is the Higgs boundary parameter,
determined by palindrome decomposition of the 13 active $f_{\rm MDL}$
neighbourhoods.

\paragraph{Frobenius-Generation Coincidence Identity (FGCI).}
The lepton ladder index $b_{L_1} = 73$ admits two \emph{independent} derivations
that agree if and only if $N_c = 3$:
(1)~the \textbf{Frobenius chain formula} $F(n) = n^4 - n^2 + 1$ evaluated at
$n = N_c = 3$ gives $81 - 9 + 1 = 73$;
(2)~the \textbf{GTE cascade formula} $G(n) = n^4 - (n^2+1)/2 - n$ at $n=3$
also gives $73$.  Both formulas coincide \emph{uniquely} at $n = N_c = 3$:
for $n \neq 3$ the two expressions differ.  This witnesses that the algebraic
structure (Frobenius group $F_{21}$) and dynamical structure (GTE cascade)
are mutually consistent precisely at the physically realised colour-charge value.
The muon ladder index $b_{L_2} = 42 = 2 \times N_c \times |Z_7|$ factors through
both generators of $F_{21}$.  The FGCI is machine-certified in Lean~4
(\lean{FrobeniusChain.lean}, 17 theorems, zero sorry, $\mathbf{A}_{\rm Lean}$;
P01~\cite{SpivackSM}).

The GTE cascade produces all SM fermion quantum numbers.  The $\mathbb{Z}_7$
winding number $W_B$ is the interaction charge; the species formula
$W_B = 4k \bmod 7$ for occupation~$k$ identifies: $k{=}1$ (charged leptons),
$k{=}4$ (up-type quarks), $k{=}5$ (down-type quarks), $k{=}7\equiv 0$ (neutrinos).
Three quark colours arise from the gen$_1$/gen$_2$ occupation ratio within a fixed
composite.  The $\mathbb{Z}_3$ color charge $Q_\chi$ is the Sylow-3 dlog of
$\mathrm{GF}(7)^*$ --- the unique Sylow-3 invariant, derived with zero SM input.
The joint topological charges $(W_B,\, Q_\chi)$ are derived from first principles
(proved $\mathbf{A/D}$, Lean-certified~\cite{SpivackGTEUnification}).

\paragraph{Temporal winding and neutrino mass origin.}
The P45 three-tape CMCA extends the winding space to
$\mathbb{Z}_7^4 = (w_x, w_y, w_z, w_\tau)$, where $w_\tau$ is the temporal
winding component of the $\tau_c$ clock field.  The neutrino sector
($k = 7 \equiv 0$, $W_B = 0$) corresponds to vacuum temporal winding $w_\tau = 0$.
The bion instanton mechanism for generating neutrino masses is ruled out at the
GTE level: the kink action parameter satisfies $\beta^2 = 49 > 8\pi$, placing
$\Phi_{\rm MDL}$ outside the bion-mass-generation regime (CatAL, Lean-certified).
Neutrino masses therefore arise from non-topological dynamics, consistent with the
seesaw structure derived in P21~\cite{SpivackNeutrino}.
The neutrino mass sum $\Sigma m_\nu = 59.4$~meV is derived from the P47 cosmological
framework via the holographic mode count applied to the neutrino sector: within the
Planck~2018 bound $< 120$~meV, predicting normal mass ordering~\cite{SpivackGTECosmology}.

Full details of the generation orbit, the Braid Atlas of topological signatures~\cite{SpivackBraidAtlas},
and the cascade derivation of the fermion spectrum are in the Standard Model
paper~\cite{SpivackSM}, the UGP dynamics paper~\cite{Spivack2026_UGPDynamics}, and
the computational universality paper~\cite{SpivackCompUniversality}.

% ====================================================================
\subsection{The Standard Model from Orbit Arithmetic}
\label{subsec:overview-sm}

The GTE cascade generates the quantitative SM parameter spectrum from the single
seed $(1,\,73,\,823)$ with zero free parameters.  Headline results
from~\cite{SpivackSM}:

\begin{itemize}[leftmargin=1.5em,itemsep=0.25em]
  \item Nine charged-fermion masses at 0.295\% RMS with zero active parameters at
    prediction time; a 1000-permutation null confirms arithmetic structure.
  \item Bare gauge couplings $g_1^2,\, g_2^2,\, g_3^2$ are Lean-certified exact
    rationals; $\alpha_s(M_Z) = 0.11822$ sits at $+0.24\sigma$ of PDG~2024, committed
    43 days before the first PDG comparison.
  \item Neutrino mass-squared ratio $\Delta m^2_{21}/\Delta m^2_{31} = 0.0294$ at
    $0.16\sigma$ from NuFIT~6.0, parameter-free.
  \item Cosmological constant from the derived model bit-length at $0.31\sigma$ from
    Planck 2018 ($H_0$ is the sole external input).
\end{itemize}

\paragraph{Electroweak mixing.}
The Weinberg angle is derived from the orbit constants in a ten-step chain, fully
machine-certified:
\[
  \sin^2\theta_W \;=\; \frac{N_{\rm gen}}{c_H} \;=\; \frac{3}{13} \;=\; 0.23077.
\]
The Wolfenstein-corrected value $384729/1664000 \approx 0.23121$ sits $-0.42\sigma$ from PDG~2022.
The two-loop SM Higgs-threshold correction with GTE-derived inputs gives
$\sin^2\theta_W = 0.23129154$, within $+0.038\sigma$ of PDG~2024
(CatB; all electroweak inputs GTE-derived except the structural identity
$1/\alpha_{\rm em}(0) = 2^{|Z_7|} + N_{\rm gen}^2 = 137$).
The $Z$-boson mass $M_Z = 91.914$
GeV ($+0.80\%$ PDG) and $W$-boson mass $M_W = 80.614$ GeV ($+0.30\%$ PDG) follow
from the same derivation chain ($\mathbf{A}_{\rm Lean}$
conditional)~\cite{SpivackWeinbergAngle}.

\paragraph{Electroweak VEV from SRRG.}
An independent first-principles derivation of the electroweak scale is
provided by the Self-Referential Renormalization Group (P27)~\cite{SpivackSRRG}.
The SRRG flow maximises a net viability functional $F[S] = R[S] - C_\Lambda[S]$
over the space of self-referential physical theories.  Its unique IR fixed
point, with linearised contraction rate $1/\varphi$ (the golden ratio), selects
the electroweak vacuum: $v_{\rm PSC} = 246.16$~GeV ($-0.024\%$ from PDG
$246.22$~GeV).  This is machine-certified:
\lean{psc\_ew\_entropy\_maximization} in \lean{GoldstoneEntropyCorrection.lean}
($\mathbf{A}_{\rm Lean}$, zero sorry, zero open axioms).
The result is independent of the Higgs mechanism; it requires only the $\mathbb{Z}_7$
structure and the PSC closure condition.
The SRRG--MDL equivalence (P27 Open Problem~9) has been substantially closed:
\lean{op9\_catal\_unconditional} proves $K_{\rm alg}(g^*) = 0$ with unique zero on
$(0,\infty)$ ($\mathbf{A}_{\rm Lean}$, zero sorry, unconditional), and
\lean{srrg\_op9\_k\_alg\_biconditional} proves the full biconditional
\textsc{IsSrrgFixedPoint} $\Leftrightarrow$ $K_{\rm alg}$ minimiser (CatAD, zero sorry,
no hypothesis).  The SRRG fixed point $g^* = 1/\varphi$ is simultaneously the MDL-optimal
coupling and the PSC-adjudication attractor: the same PSC axioms that select $N_{\rm gen}=3$
also drive the SRRG flow to $g^*$ (CatAD, srrg-lean).

\paragraph{SU(2)$_L$ gauging.}
The within-tape weak-isospin doublet $(\Phi_+,\Phi_-)$ has $SU(2)_L$-invariant
potential $V(|\Psi|)$ since $|U\Psi|=|\Psi|$.  The MDL principle forces minimal
local gauging with zero extra specification cost ($\CatAL$, zero named axioms):
\lean{phimdl\_potential\_su2l\_invariant} certifies $L_2$ norm invariance under
$SU(2)$ rotation; \lean{su2l\_covariant\_derivative\_minimal} certifies the
covariant derivative as MDL-minimal; \lean{su2l\_wpm\_generator\_algebra}
certifies that $W^\pm$ satisfy the $SU(2)$ Lie algebra.  The master bundle
\lean{su2l\_l2\_from\_phimdl\_potential\_catad} packages the full gauging chain
(P46~\cite{SpivackGTEPolynomialUFT}, \texttt{GaugeMDL.lean}).
The quantitative $V{-}A$ chiral structure of the charged current
($g_R = 0$, $g_A/g_V = -1$) now follows from the opposite tape windings
$w_L = -w_R$ of the CMCA at CatAD (\lean{va\_coupling\_exact\_from\_tape\_topology}).

\paragraph{CKM mixing.}
All four Wolfenstein parameters of the CKM quark mixing matrix are derived from the
GTE orbit constants with zero free parameters.  The leading Wolfenstein parameter
\[
  \lambda \;=\; \frac{N_{\rm gen}^2}{2^{N_{\rm gen}} \cdot N_{\rm fam}} \;=\;
  \frac{9}{40} \;=\; 0.225,
\]
agrees with the PDG value to $0.0\sigma$ (Lean-certified)~\cite{SpivackCKM}.
The hierarchy parameters $A$, $\bar\rho$, $\bar\eta$ follow from the quark-sector
orbit indices, with the measured CP-violating phase~$\delta_{\rm CKM}$ reproduced
without input.

\paragraph{Fine structure constant.}
The electromagnetic coupling constant
$\alpha_{\rm EM}^{-1} = 137$ is derived via two independent routes: from the GTE
cascade chain (Lean: \lean{alpha\_inv\_master}, $\mathbf{A}_{\rm Lean}$) and from
$F_{21}$ Berry phase holonomy on a single $\Phi_{\rm MDL}$ kink substrate (ROBUST,
$\mathbf{A}_{\rm Lean}$, eight independent audit passes)~\cite{SpivackGTEUnification}.

\paragraph{Lepton mass scale.}
The $\mathbb{Z}_7$ potential coefficient $c_V = 1/49 = 1/N_{Z_7}^2$ determines
the tau-lepton Yukawa coupling exactly:
\[
  y_\tau \;=\; \frac{c_V}{N_{\rm mod2}} \;=\; \frac{1}{2 \times 7^2} \;=\; \frac{1}{98},
\]
where $N_{\rm mod2} = 2$ is the $\mathbb{Z}_2$ parity factor.  All three
charged-lepton masses follow from $v_H$ alone:
\[
  m_\tau \;=\; \frac{v_H}{98\sqrt{2}} \;=\; 1776.57~\text{MeV} \quad (0.016\%).
\]
Via the Koide relation ($\theta = 2/9$ from $N_c = 3$; Lean-certified:
\lean{koide\_angle\_from\_N\_c\_pure}), $m_\mu = 105.6$~MeV ($0.022\%$) and
$m_e = 0.511$~MeV ($0.023\%$) follow with zero additional SM inputs.
The sole dimensional anchor is $G_F$ (fixing $v_H = (\sqrt{2}\,G_F)^{-1/2}$);
no fermion mass anchor is required.

The Koide relation admits the reformulation $\mathrm{CV}(\sqrt{m}) = 1$: the
coefficient of variation of the square-root masses equals unity, characterising
the equal-irrep-norm (maximally-dispersed positive spectrum) condition.  This form
is Lean-certified (CatAL): $Q = (1 + \mathrm{CV}^2)/3 = 2/3$ iff $\mathrm{CV} = 1$.
The kink--Higgs coupling derived from $c_V$ is
$g_{hKK} = 4/7^4 = M_{\rm kink}/(v_H/\sqrt{2})$,
with the identity $y_\tau = c_V/N_{\rm mod2}$ algebraically exact (CatA).
The cyclotomic structure underlying the Koide relation is established in P18~\cite{Spivack2026_Koide}.

% ====================================================================
\subsection{QCD and the Strong Force from \texorpdfstring{$F_{21}$}{F21}}
\label{subsec:overview-qcd}

The substrate's internal discrete symmetry group is identified as
$F_{21} = \mathbb{Z}_7 \rtimes \mathbb{Z}_3 = \Sigma(21) \subset \mathrm{SU}(3)$ ---
the unique non-abelian group of order~21 --- forced by MDL minimality (Lean-certified,
$\mathbf{A}_{\rm Lean}$)~\cite{SpivackCompUniversality,SpivackGTEUnification}.
Its abelianisation $F_{21}^{\rm ab} = \mathbb{Z}_3$ is the color gauge symmetry; its
non-abelian structure at the group level provides exact QCD Casimir invariants and
structure constants without introducing non-abelian gauge bosons in the low-energy
Lagrangian.

From $F_{21}$ the following QCD results are derived:

\begin{itemize}[leftmargin=1.5em,itemsep=0.25em]
  \item \textbf{Asymptotic freedom:} $b_0 = (11N_c - 2N_f)/3 = 7 = |{F_{21}}|/3$,
    with $N_c = 3$ and $N_f = 2N_{\rm gen} = 6$ derived from the GTE structure.
    $b_0 = |\mathbb{Z}_7|$ is Lean-certified
    (\lean{b0\_eq\_z7\_order}, \lean{f21\_substrate\_asymptotic\_freedom},
    $\mathbf{A}_{\rm Lean}$); the Frobenius prime identity
    $|\mathbb{Z}_7| = |\mathbb{Z}_3|^2 - |\mathbb{Z}_3| + 1$ is separately
    certified~\cite{SpivackGTEQCDStructure}.
  \item \textbf{One-loop $\alpha_s$:} $\alpha_s(M_Z) = 0.1319$ (one-loop);
    $\alpha_s = 0.1201$ at two loops with $N_f{=}5$ threshold --- within 2\% of
    PDG~\cite{SpivackGTEQCDStructure}.
  \item \textbf{Colour confinement:} No PSC-admissible single-quark state exists
    (Lean: \lean{no\_psc\_admissible\_single\_quark}, $\mathbf{A}_{\rm Lean}$);
    every physical state is PSC-admissible, confined, massive, and anomaly-free
    (\lean{unified\_physical\_exclusion}, \lean{confined\_massive\_color\_singlet},
    $\mathbf{A}_{\rm Lean}$)~\cite{SpivackEmergentGravity}.
  \item \textbf{Strong CP problem:} $\theta_{\rm QCD} = 0$ exactly, derived from
    three independent $F_{21}$ group-theory arguments:
    $\pi_3(F_{21})=0$, element-wise $\det = 1$, and $\sum{\rm Im}(\chi_3)=0$.
    All Lean-certified with zero sorry; no axion mechanism is
    required~\cite{SpivackGTEQCDStructure}.
  \item \textbf{Hadron spectroscopy:} The meson pseudoscalar nonet (9 states),
    baryon octet (8), baryon decuplet (10), and vector meson nonet (9) are all
    reproduced from GTE kink composites.  Casimir invariants $C_F = 4/3$,
    $C_A = 3$, $T_F = 1/2$, and all nine structure constants $f^{abc}$ are
    reproduced at machine precision from $F_{21}$ generators.  The $\eta$--$\eta'$
    mixing angle $\theta_P = -13.08^\circ \pm 3.74^\circ$ lies in the PDG range
    $[-14.3^\circ, -10.7^\circ]$ with zero PDG inputs.  All six quark masses are
    reproduced within 7\% of PDG current values~\cite{SpivackGTEQCDStructure}.
  \item \textbf{All 8 colour gluons} emerge from $F_{21}$ Berry holonomy on a single
    $\Phi_{\rm MDL}$ kink substrate (PROVISIONAL-STRONG)~\cite{SpivackGTEUnification}.
  \item \textbf{SU(3) UV completion ($\mathbf{A/D}$):}
    Below $\Lambda_{\rm GTE} \approx 2$~GeV, the $F_{21}$ abelian gauge effective
    theory is the correct description.
    Above $\Lambda_{\rm GTE}$, $F_{21} \hookrightarrow \mathrm{SU}(3)$ faithfully
    ($\mathbf{A}_{\rm Lean}$: \lean{F21SU3Embedding.lean}); the Burnside density
    theorem establishes $\overline{\mathrm{span}}_\mathbb{C}\,\rho(F_{21}) =
    M_3(\mathbb{C})$, and $\Phi_{\rm MDL}$ scalar fluctuations fill the
    $\mathrm{SU}(3)/F_{21}$ coset, completing the UV theory as full SU(3)
    Yang--Mills.
    The na\"ive Wilsonian-flow approach from $F_{21}$ to $\mathrm{SU}(3)$ is
    falsified: the discrete $F_{21}$ lattice does not continuously flow to
    $\mathrm{SU}(3)$ in the IR; the UV completion is via Burnside density, not
    RG flow~\cite{SpivackGTEUnification}.
  \item \textbf{Predicted second Cartan gauge field $A'_\mu$:}
    The $F_{21} \subset \mathrm{SU}(3)$ embedding has rank~2 (two Cartan
    generators $\lambda_3$, $\lambda_8$).  The current GTE Lagrangian provides
    only the first Cartan field $A_\mu$ ($H_A = (-T^3 + \sqrt{3}\,T^8)/2$).
    A second Cartan field $A'_\mu$ ($H_{A'} = (\sqrt{3}\,T^3 + T^8)/2$) is
    predicted but currently absent from the Lagrangian.
    Its coupling $e' = e$ is forced by Killing-form equality
    $\mathrm{tr}(H_A^2) = \mathrm{tr}(H_{A'}^2) = 1/2$; no new free parameters
    are introduced ($\mathbf{A}_{\rm Lean}$).
    In QCD, $\lambda_3$ and $\lambda_8$ are both present; GTE is predicting the
    correct colour-sector structure.
\end{itemize}

The mass gap is established at the lattice level (ROBUST, $\mathbf{A}_{\rm Lean}$
lower bound) with continuum-limit scaling confirmed~\cite{SpivackEmergentGravity}.
A non-perturbative string tension computation on the three-dimensional $G_{13}$
deformation lattice confirms $b_0 = 7 = |\mathbb{Z}_7|$ as the dominant RG
coefficient; the non-perturbative contribution dominates over the one-loop term,
and the quantisation factor $f_{\rm quant}$ is PROVISIONAL CatA~\cite{SpivackGTEQCDStructure}.

The algebraic identity of $F_{21}$ within the broader PSL(2,7) structure has been
machine-certified in P52~\cite{SpivackPSL27Unification}: $F_{21}$ is the Borel
stabilizer of $\infty$ in the unique simple group $\mathrm{PSL}(2,7)$ of order~168,
which simultaneously acts as the automorphism group of the Fano plane (with $F_{21}$
as its flag-transitive subgroup) and as the automorphism group of the Klein quartic
surface of genus $N_{\rm gen} = 3$.
The Eisenstein functor $\mathbb{Z}[\omega]$ unifies $F_{21}$ (via the split prime~7)
and $A_4 = V_4 \rtimes \mathbb{Z}_3$ (via the inert prime~2) from the same ring.
All four results are machine-certified with zero sorry ($\CatAL$).

% ====================================================================
\subsection{Computational Universality: Cook's Theorem and Finite NAND Completeness}
\label{subsec:overview-cup}

Rule~110 is universally regarded as the canonical example of a computationally
universal elementary cellular automaton, established by Cook's proof
(2004)~\cite{Cook2004}.  The conditional operational readback core of Cook's
construction is machine-certified in Lean~4
(\lean{cook\_operational\_stage3\_tm\_microstep\_readback}, zero sorry modulo five named
classical Cook collision-analysis bridge axioms)~\cite{SpivackCookTheorem} --- a
certificate for already-supplied compilations, not itself a self-contained
Turing-universality theorem.  Cook's classical theorem remains the load-bearing
result for Rule~110's Turing universality.

Independently of Cook's construction, the UGP Physics programme establishes an exact
polynomial representation of the Rule~110 update rule over the Galois field
$\mathrm{GF}(7) = \mathbb{Z}/7\mathbb{Z}$:
\[
  p(L,\,C,\,R) \;=\; C + R - C{\cdot}R - L{\cdot}C{\cdot}R,
\]
where $L, C, R \in \{0,1\} \subset \mathrm{GF}(7)$.  When $C = 1$, this simplifies
to $p(L,1,R) = {\rm NAND}(L,R)$.  Since NAND is functionally complete (Sheffer
1913), Rule~110 at center cell $C=1$ realizes any finite two-input Boolean function
algebraically, without reference to Cook's construction --- a finite functional-
completeness result about a fixed input size, distinct from and not sufficient for
Turing universality, which requires simulation on an unbounded tape for an unbounded
number of steps.  The polynomial and NAND identity are Lean-certified with zero sorry
(Lean: \lean{rule110\_z7\_poly\_rep}, \lean{rule110\_center1\_is\_nand})~\cite{SpivackGF7Universality}.
$\Phi_{\rm MDL}$'s Turing universality is established via the Cook route composed with
the proved stepwise $\Phi_{\rm MDL}$--Rule~110 simulation
(\lean{phimdl\_turing\_universal}, $\mathbf{A}_{\rm Lean}$ conditional; zero sorry;
one named axiom: \lean{cook\_rule110\_simulates\_computable}); the UWCA register file
provides a second, independent Turing-universality route via Minsky counter-machine
simulation, conditional on one named classical axiom
(\lean{minsky\_counter\_machine\_turing\_complete\_1967}).

A further consequence is the \textbf{Algebraic Descent Theorem} (Lean-certified,
$\mathbf{A}_{\rm Lean}$): every algebraic and topological property of $\Phi_{\rm MDL}$
that depends only on the $F_{21}$ internal symmetry holds for $f_{\rm MDL}$ at any
lattice resolution $M \geq 1$.  Three generations, colour confinement, asymptotic
freedom, and strong CP are all resolution-independent --- proved at the algebraic level
and therefore simultaneously valid for the discrete algebraic certificate and the
continuum physical field theory~\cite{SpivackGTEUnification}.

% ====================================================================
\subsection{Quantum Mechanics from the \texorpdfstring{$\Phi_{\rm MDL}$}{Phi\_MDL} Field}
\label{subsec:overview-qm}

Quantum mechanics in the GTE framework is not postulated; it is derived from the
structure of the $\Phi_{\rm MDL}$ field and the $[D]$-selection
mechanism~\cite{SpivackQMFromR110,SpivackPhiMDLField}.

\paragraph{Two levels of description.}
At the Level-1 algebraic beable level (the CMCA discrete certificate), orbit states
are always definite: no superposition exists at this level.  At the Level-2 continuum,
the $\Phi_{\rm MDL}$ field admits superpositions
$\sum_k c_k\lvert{\rm kink}_k\rangle$ over its $\mathbb{Z}_7$ winding sectors.
Quantum mechanics emerges entirely at Level~2; the discrete certificate provides
the algebraic labelling of the winding sectors.

The selection mechanism rests on the GTE-Möbius architecture framework~\cite{SpivackGTEMobius}: the abstract triple $(A, e, [D])$ where $A$ is the Z$_7$ arithmetic carrier, $e$ is the GF(7)/Cook self-encoding map, and $[D]$ is the class of PSC-consistent coherence measures characterised by five constraints (D1--D5: non-negativity, PSC invariance, non-computability, selection completeness, Born consistency). The abstract carrier $A$ has its concrete physical realisation as the PSC-admissible kink sector of $\Phi_{\rm MDL}$; the five constraints are established in P34~\cite{SpivackGTEMobius} and serve as the axiomatic foundation of the quantum field-theoretic results in P42--P43. Two functions implement quantum mechanics on this substrate:

\begin{enumerate}[leftmargin=2em,itemsep=0.15em]
  \item The \textbf{Born rule} $\mathcal{B}$: maps amplitudes to
    probabilities $P(k) = |c_k|^2$.  Derived unconditionally and Lean-certified
    with zero custom axioms: \lean{born\_rule\_unconditional}
    ($\mathbf{A}_{\rm Lean}$)~\cite{SpivackPhiMDLField}.  An independent derivation
    via the Page--Wootters (PW) mechanism has been established at CatAD (P45): the
    $\tau_c$ clock, as a diagonal-Hamiltonian (Salecker--Wigner--Peres semiclassical)
    clock, satisfies all PW prerequisites, and the timeless universe state
    $(1/\sqrt{T})\sum_\tau|\tau\rangle_{\rm clock}\otimes U(\tau)|\psi_0\rangle_{\rm sys}$
    yields $P(k|\tau) = |\langle k|U(\tau)|\psi_0\rangle|^2$ at every $\tau$,
    independently of whether $\tau_c$ is in a classical or quantum superposition.
    The two derivations are mutually consistent and reinforce each other.
  \item \textbf{Transputation} $P^\top$: selects the physical state
    $\rho^* = \mathrm{argmin}_\rho\, D(\rho\,|\,w)$ from record class~$w$ via the
    $[D]$-functional ($\mathbf{A/D}$; Lean: \lean{TransputationStateSelector.lean})~\cite{SpivackSDS_PR0}.
\end{enumerate}

\paragraph{Objective reduction.}
$[D]$-selection is objective (no conscious observer required), determinate (unique
minimum from D4), and non-computable (D3 places it outside Turing-computable
selection).  The non-computability is now characterised precisely: the adjudication
function is exactly $\Delta_2^0$-Turing-complete (the Gold--Putnam
``trial-and-error'' stratum, degree $0'$) --- record truth on the diagonal fragment
is $\Sigma_1^0$-complete, while the adjudicator is limit-computable, so quantum
measurement operates precisely at the halting degree, never requiring oracle access
above the first Turing jump (P51~\cite{SpivackPolynomialTransputation}).
This dissolves the measurement problem --- reframing it in terms of PSC-forced transputation
(conditional on the D4 universality assumption) --- without many-worlds, pilot waves,
or retrocausation.  EPR correlations are semantic: they exist in the $[D]$-observable
layer as constraints from the shared record structure established at interaction time;
the beable layer is always local and deterministic.  GTE escapes Bell's theorem
because the transputation mechanism~D3 is non-computable, placing it outside Bell's
local-computable-hidden-variable assumptions.  Delayed-choice and quantum eraser
experiments are resolved by the same two-level structure: the CA beable trajectory is
always definite at Level~1; whether an interference pattern or which-path statistics
is observed depends only on which records~$w$ are retained in the $[D]$-description,
with no retroactive modification of any beable-level
trajectory~\cite{SpivackPhiMDLField}.

\paragraph{Thermal state.}
Prior to any $[D]$-selection event, the physical state is the Hamiltonian thermal
ensemble on PSC-admissible kink sectors $\{0,2,3,4,6\}$: vacuum-dominant at
cosmological temperatures (Lean-certified: \lean{PhiMDLThermalState.lean},
$\mathbf{A}_{\rm Lean}$ conditional).

% ====================================================================
\subsection{Special Relativity and the Three-Layer Chiral Minkowski CA}
\label{subsec:overview-sr}

Lorentz invariance is not postulated in the GTE framework.  It emerges at the
observer level from a specific two-level structure~\cite{SpivackCMCA}.

\paragraph{The CMCA.}
The \textbf{Three-Layer Chiral Minkowski CA (CMCA)} is the MDL-minimal Level-1
discrete construction realising all five required structural features: Turing
universality, $\mathbb{Z}_7$ orbit, $V{-}A$ weak-current chirality, observer-level
Lorentz invariance, and dynamics-level special relativity.  It consists of an outer
chiral CA pair (Rule~110 + Rule~124) plus an inner $\tau_c$ Rule-110 clock per cell.
The MDL bit-count $K_{\rm CA} = 19$ bits saturates the construction-class lower bound
(Lean: \lean{CMCAMDLMinimality.lean}, $\mathbf{A}_{\rm Lean}$ conditional).

\paragraph{Emergence of Lorentz symmetry.}
At the CA beable level, the Rule~110 lattice carries a preferred rest frame; Lorentz
symmetry is broken by $\varepsilon_0(M) = \pi^2/(3M^2)$ at the lattice scale.
PSC Presentation Invariance (PI) forces the $[D]$-selection mechanism to be
Lorentz-equivariant (Lean: \lean{PSCPILorentzMain.lean}, $\mathbf{A}_{\rm Lean}$):
no $[D]$-selected observable can distinguish inertial frames.  Lorentz symmetry is a
property of the selection mechanism, not of the underlying beables.  At Planck-scale
lattice spacing $a = \ell_P$, the lattice correction $\varepsilon_0 \approx 10^{-62}$
is observationally indistinguishable from zero.

\paragraph{Continuum limit.}
The continuum limit of the CMCA is identified with $\Phi_{\rm MDL}$
($\mathbf{A}_{\rm Lean}$ conditional): \lean{CMCAContinuumLimit.lean}.  The
$\Phi_{\rm MDL}$ field satisfies exact Lorentz invariance $\omega^2 = k^2 + m^2$,
machine-certified in Lean~4 (\lean{poincare\_invariance\_of\_kg},
$\mathbf{A}_{\rm Lean}$)~\cite{SpivackCMCA,SpivackPhiMDLField}.

% ====================================================================
\subsection{Gravity from the \texorpdfstring{$\Phi_{\rm MDL}$}{Phi\_MDL} Field}
\label{subsec:overview-gravity}

Gravity emerges from the $\Phi_{\rm MDL}$ field structure in three interconnected
results~\cite{SpivackPhiMDLGravity,SpivackCompleteness}:

\paragraph{Two routes to emergent gravity.}
Gravity emerges from the GTE framework via two complementary routes.
The \textbf{P36 route} (CMCA/discrete level) derives geodesic structure from the
deviation-based Ollivier--Ricci curvature of the Rule~110 causal graph:
the discrete vacuum Einstein equation $G_{\mu\nu}=0$ is machine-certified
(\lean{gte\_spacetime\_dimension}, $\mathbf{A}_{\rm Lean}$); the spectral
dimension $d_s = 4$ is separately machine-certified
(\lean{causal\_graph\_spectral\_dim\_thermodynamic\_limit}, $\mathbf{A}_{\rm Lean}$).
An MDL--Lovelock structural correspondence ($\mathbf{A/D}$) establishes that
Einsteinian gravity emerges as the unique diffeomorphism-invariant effective
theory compatible with the GTE substrate~\cite{SpivackEmergentGravity}.
This constitutes the \emph{structural/geometric certificate}: gravity is
forced by the information structure.
The 1D discrete gravity simulation (P36) yields a negative or null result for the inter-kink
force in 1+1D, as expected: a 1+1D CMCA tape does not have three independent spatial directions
and therefore cannot produce an inverse-square-law in the familiar sense; this is not a failure
but a structural feature of the 1+1D algebraic certificate level.
The \textbf{P45 three-tape route} (3+1D continuum level) closes this gap.
With three independent CMCA tapes sharing a common clock $\tau_c^{\rm out}$ (DPP),
the PMDL Poisson equation $\nabla^2\Phi = G_{\rm eff}\,p(w_x,w_y,w_z)$ yields
$F\propto b^{-2}$ in the far-field limit, machine-certified and confirmed analytically
($\mathbf{A/D}$, CatAD; see \S\ref{subsec:overview-substrate}).
The two routes are not competing: P36 shows \emph{why} gravity is forced (structural/geometric);
P38/P45 show \emph{how} it manifests quantitatively in the continuum EFT.

\paragraph{Stress-energy tensor and Einstein equations.}
The stress-energy tensor
$T_{\mu\nu} = \partial_\mu\Phi\,\partial_\nu\Phi - \eta_{\mu\nu}\mathcal{L}$ is
derived analytically from the $\Phi_{\rm MDL}$ Lagrangian and Lean-certified
($\mathbf{A/D}$): symmetry, vacuum-vanishing, and BPS pressure-free condition
$T_{11} = 0$ are all proved in \lean{StressEnergyTensor.lean}.  The kink mass
$\int T_{00} = 290.10$~MeV is verified numerically to $0.01\%$; energy-momentum
conservation holds to $6 \times 10^{-12}$.

The Einstein equations $G_{\mu\nu} = 8\pi G\,T_{\mu\nu}[\Phi_{\rm MDL}]$ are derived
via MDL-Lovelock (the unique dimensionally correct generally covariant action in
$d{=}4$), minimal coupling, and variational calculus ($\mathbf{A/D}$).  The
classical cosmological constant vanishes exactly: the $\mathbb{Z}_7$-symmetric
potential $V(\phi) = m^2(1-\cos 7\phi)/49$ has $\langle V \rangle = 0$ in the PSC
vacuum by exact cancellation, Lean-certified.  The $\mathbb{Z}_7$ global shift
symmetry $\phi\to\phi+2\pi/7$ is furthermore \emph{anomaly-free}: the scalar
path-integral measure $\prod_x d\phi(x)$ is trivially translation-invariant (no
Jacobian anomaly, unlike fermionic determinants), so all 7 vacuum sectors
remain quantum-degenerate to all loop orders \emph{in the pure-$\varphi$ sector}
--- the degeneracy theorem is hypothesis-scoped to $\mathbb{Z}_7$-periodic
observables.  Four zero-sorry Lean theorems
certify this: \lean{z7\_shift\_measure\_preserving}, \lean{z7\_jacobian\_eq\_one},
\lean{lebesgue\_shift\_invariant}, and \lean{z7\_vacuum\_sectors\_equiprobable}
($\mathbf{A}_{\rm Lean}$; P47/P43).  Quantum vacuum degeneracy of the pure-$\varphi$
sector is therefore a structural theorem, not an assumption.  In the full Lagrangian,
the canonical coupling $V_{\rm coupling} = \varepsilon|\phi|^2(D_\mu\chi)^2$ is not
$\mathbb{Z}_7$-periodic and distinguishes the seven vacua through the $\chi$-sector
kinetic normalization; this controlled breaking of the vacuum-label degeneracy
supplies the vacuum-selection bias that resolves the cosmological domain-wall
problem (\S\ref{subsec:overview-complete};
P42, P47~\cite{SpivackPhiMDLField,SpivackGTECosmology}).

\paragraph{The gravitational $D^2$ theorem.}
The factor $D^2 = (N_{\rm gen}+1)^2 = 4^2 = 16$ unifies the gravitational and
cosmological sector of GTE.  With $D = N_{\rm gen}+1 = 4$ derived from the DPP
($\mathbf{A}_{\rm Lean}$, \lean{cmca\_tensor\_product\_gives\_31d\_minkowski}), every
area-type gravitational normalization carries exactly one factor of $D^2$:
\begin{itemize}[noitemsep,topsep=2pt]
  \item CC epoch count: $D^2 N_{\rm fam}^{N_{\rm gen}}/N_{\rm gen} = 2000/3$ (CatAD)
  \item CC magnitude: $\rho_\Lambda = N_{\rm spatial}^2/(D^2|\mathbb{Z}_7|)\,M_{\rm Pl}^2 H_0^2$ (CatAD)
  \item Gorard coefficient: $C_{\rm Gorard} = N_{\rm spatial}/(2D^2) = 3/32$
    ($\mathbf{A}_{\rm Lean}$, \lean{c\_gorard\_eq\_n\_spatial\_over\_two\_dsq})
  \item Einstein--Hilbert normalization: $1/(D^2\pi G_N)$ (standard GR; $D=4$ from DPP)
\end{itemize}
$D^2$ counts ordered pairs of orbit discriminants in $D$-dimensional spacetime.
It does \emph{not} appear in the holographic sector ($\Omega_\Lambda^{\rm holo}$,
$F_{21}$, $\sin^2\theta_W$).  In particular, the Gorard bridge coefficient
$C_{\rm Gorard} = N_{\rm spatial}/(2D^2) = 3/32$ is an \emph{exact} algebraic
identity proved in Lean~4 with zero sorry --- not an approximation.

\paragraph{Newton's constant from the Planck/tau hierarchy.}
Newton's constant is derived from the ratio:
\[
  \frac{M_{\rm Pl}}{m_\tau} \;=\; \frac{F_{21}^{10} \times |\mathbb{Z}_7|^7}{2},
\]
where $n = 10$ is the GTE orbit decomposition count on $\mathbb{Z}_7^3$
(all-distinct orbits equal the PSC count $(N_{\rm gen}-1)\times N_{\rm cells} = 10$),
$b_0 = |\mathbb{Z}_7| = 7$ is derived from $F_{21}$ group theory, and the factor
$1/2 = |\mathbb{Z}_2|^{-1}$ reflects the chirality of the Rule-110/Rule-124 pair.
The resulting prediction $M_{\rm Pl}^{\rm GTE} = 1.2204 \times 10^{22}$~MeV agrees
with the PDG value to 0.040\% ($\mathbf{A/D}$, Lean-certified conditional:
\lean{mpl\_mtau\_formula\_catal\_conditional}, \lean{b0\_eq\_z7\_order},
\lean{FrobeniusPrimeIdentity.lean})~\cite{SpivackCompleteness}.

\paragraph{Geodesic theorem and spectral dimension.}
Free particles follow paths of minimum computational action through the
$\tau_c$-metric curved by matter.  The geodesic theorem is Lean-certified to zero
axioms in three forms: timelike worldlines
(\lean{psc\_orbit\_is\_curvature\_geodesic}), flat spatial paths
(\lean{geodesic\_theorem\_spatial\_general}), and the Ehrenfest reduction
(\lean{dweight\_centroid\_tracks\_geodesic}), all $\mathbf{A}_{\rm Lean}$, zero
sorry~\cite{SpivackCompleteness}.  The spectral dimension $d_s = 4$ in the
thermodynamic limit is separately derived~\cite{SpivackEmergentGravity}.

\paragraph{Quantum gravity and holography.}
The curved-background Lagrangian $\mathcal{L}[\Phi_{\rm MDL};\,g_{\mu\nu}]$ is uniquely
determined: MDL minimality forces $\xi = 0$
(\lean{minimal\_coupling\_is\_mdl\_minimal}, $\mathbf{A}_{\rm Lean}$), yielding the
minimal-coupling form
$\mathcal{L} = \tfrac{1}{2}g^{\mu\nu}\partial_\mu\Phi\,\partial_\nu\Phi - V(\Phi)$
with no additional curvature terms ($\mathbf{A/D}$).  The full quantum gravity sector
derives the Graviton Fock space
$\mathcal{H}_{\rm phys} = \mathcal{H}_{\Phi_{\rm MDL}} \otimes \mathcal{H}_{\rm grav}$
and Bekenstein--Hawking entropy $S_{\rm BH} = A M_{\rm Pl}^2/4$ by two independent
routes (Wald entropy, with $\xi = 0$; and domain wall).  The holographic entropy
formula
\[
  S(A) = \frac{\mathrm{Area}[\Gamma_{\rm min}]}{4G}
\]
is established by both routes; the Reed--Solomon quantum error correction route is
additionally Lean-certified
(\lean{gte\_orbit\_vectors\_are\_rs\_codewords}, $\mathbf{A}_{\rm Lean}$, zero sorry).
The Hawking temperature $T_H = M_{\rm Pl}^2/(8\pi M_{\rm BH})$ is confirmed unchanged
by the field mass $m_\varphi = m_\tau$ (the potential vanishes at the horizon via the
Schwarzschild redshift factor); the critical mass
$M_{\rm crit} = M_{\rm Pl}^2/(8\pi m_\tau) = 3.34 \times 10^{39}$~MeV marks the
onset of exponential emission suppression ($\mathbf{A/D}$).  BH evaporation is unitary:
GKSL dynamics admit a Stinespring dilation confirming information
preservation~\cite{SpivackBHUnitarity}.
The Planck density state count $7^8 > e^{4\pi}$ is Lean-certified
(\lean{planck\_density\_bound\_via\_lifting}, $\mathbf{A}_{\rm Lean}$, zero sorry),
confirming that the CMCA lattice governs $E \geq M_{\rm Pl}$; singularities are
resolved at $\rho_{\rm max} = M_{\rm Pl}^4$, and the partition function $Z_{\rm GTE}$
is UV-finite by $\mathbb{Z}_7$-compactness ($\mathbf{A/D}$).
The MDL-minimal cosmological initial state is flat ($k = 0$), spatially uniform, and
kinetically dominated at $\dot\Phi_0 = M_{\rm Pl}$: flatness, uniformity, and the
absence of domain walls all follow from MDL complexity minimisation, dissolving the
need for inflation without introducing new fields ($\mathbf{A/D}$).  A GTE-specific
Friedmann correction produces a non-singular bounce at $\rho = \rho_{\rm Pl}$
(confirmed numerically, $\mathbf{A}$).  The quantum cosmological constant
($\sim 10^{45}$ vacuum-energy hierarchy) remains open; the classical $\Lambda = 0$ is
Lean-certified~\cite{SpivackCompleteness}.

P44~\cite{SpivackQGR} completes the perturbative quantum gravity sector.
UV finiteness on arbitrary smooth curved backgrounds is established via the
DeWitt--Schwinger heat-kernel expansion: all curved-background UV contributions reduce
to finite Planck-scale renormalizations ($C_i \sim 41.76$), leaving no residual UV
problem.  The \textbf{GTE Holographic Encoding Theorem} relates five descriptions
--- Lagrangian, Reed--Solomon $[5,3,3]_7$ algebraic code, holographic/RT, MDL-information-theoretic,
and quantum error correction.  Twelve of the twenty directed implications, linking
the Lagrangian, Reed--Solomon, MDL, and QEC descriptions to one another, are
established by derivation with no gravitational input; the remaining eight, connecting
any of those four to the holographic/RT description, hold conditional on the standard
Bekenstein-Hawking coefficient $1/4G$ taken as an external premise, not derived from
GTE structure.  The MDL cosmological initial state requires only $\log_2 3
\approx 1.585$ bits of initial data.  At leading order the primordial spectral index
satisfies $n_s = 1$ exactly; a seven-step conditional derivation chain yields
\[
  n_s \;=\; 1 - \frac{\ln 2}{2\pi^2} \;=\; 0.96488,
\]
matching Planck~2018 at $-0.005\sigma$ when the EFE-bridge step is established
(5 of 7 steps closed).  The tensor-to-scalar ratio satisfies
\[
  r \;=\; 0
\]
exactly, constituting the primary falsifiable prediction for the LiteBIRD satellite
($\sigma_r \approx 0.001$).

% ====================================================================
\subsection{The Complete \texorpdfstring{$\Phi_{\rm MDL}$}{Phi\_MDL} Framework}
\label{subsec:overview-complete}

Papers~P42 and~P43~\cite{SpivackPhiMDLField,SpivackCompleteness} constitute the
capstone of the 1+1D programme.  P42 establishes four principal results for
$\Phi_{\rm MDL}$ as a field theory: exact Poincar\'{e} invariance
($\mathbf{A}_{\rm Lean}$), the Born rule from field amplitudes
(\hspace{0pt}$\mathbf{A}_{\rm Lean}$), the uniqueness of $\Phi_{\rm MDL}$ among all
$\mathbb{Z}_N \times \mathbb{Z}_M$ sine-Gordon fields satisfying PSC-axiom-selected
constraints (MDL-proved), and the no-CA-replica theorem
(\lean{no\_finite\_ca\_exact\_lorentz\_replica} + \lean{phimdl\_is\_unique\_exact\_lorentz\_model},
$\mathbf{A}_{\rm Lean}$).

P43 closes the framework with three capstone results:

\begin{enumerate}[leftmargin=2em,itemsep=0.2em]
  \item \textbf{Algebraic Necessity Theorem} ($\mathbf{A}_{\rm Lean}$,
    \lean{AlgebraicNecessityTheorem.lean}): MDL minimality forces the unique
    $\mathbb{Z}_7 \times \mathbb{Z}_3$ structure over all competitors.
  \item \textbf{Complete Lagrangian derivation:} the $\Phi_{\rm MDL}$ Lagrangian
    is fully derived from $F_{21}$ symmetry with zero free parameters.  The
    coupling form $V_{\rm coupling} = \varepsilon|\phi|^2(D_\mu\chi)^2$ and coupling
    constant $\varepsilon = N_7/N_3^2 = 7/9 \approx 0.778$ (in the BPS range
    $[0.444, 0.800]$) are both Lean-certified ($\mathbf{A}_{\rm Lean}$).  The field
    mass $m_\varphi = m_\tau = 1776.86$~MeV follows from the $F_{21}$ leptonic-sector
    self-consistency condition.  The kink--Higgs coupling
    $g_{hKK} = 4/7^4 = M_{\rm kink}/(v_H/\sqrt{2})$ and the tau-lepton Yukawa
    $y_\tau = c_V/N_{\rm mod2} = 1/98$ are both algebraically exact consequences
    of $c_V = 1/49$ (CatA).
  \item \textbf{Newton's constant at high precision} ($\mathbf{A/D}$, 0.040\%,
    upgraded in P43).
\end{enumerate}

The complete two-level algebraic structure is established by the interplay of two
theorems: the \textbf{Algebraic Lifting Theorem} (Lean: \lean{algebraic\_lifting\_theorem},
$\mathbf{A}_{\rm Lean}$) carries algebraic results from discrete to physical scale;
the \textbf{Algebraic Descent Theorem} (Lean: \lean{AlgebraicDescentTheorem.lean},
$\mathbf{A}_{\rm Lean}$) shows that all $F_{21}$-dependent properties hold at every
lattice resolution.  Results proved for the continuous $\Phi_{\rm MDL}$ are
simultaneously results about the discrete $f_{\rm MDL}$ CA --- and vice versa ---
with no additional proof required.

Papers P44--P47~\cite{SpivackQGR,SpivackThreeTapeCMCA,SpivackGTEPolynomialUFT,SpivackGTECosmology}
constitute the 3+1D extension, quantum-gravity completion, and cosmological predictions.
P44 (quantum gravity completeness) establishes UV finiteness on curved backgrounds
and derives the primary cosmological falsifiable predictions.
P45 (three-tape CMCA) proves the Dimensional Protocol Principle that identifies the
minimal causal structure producing 3+1D spacetime from 1+1D building blocks, derives
the SM particle spectrum from three-tape winding arithmetic, and establishes all
structural features of 3+1D physics (SR, gravity, QM, confinement, Bell nonlocality).
P46 (polynomial UFT) establishes MDL uniqueness of $p(L,C,R)$ and derives its
five physical roles, unifying the algebraic and physical perspectives.
P47 (cosmological predictions) derives the dark-energy fraction via two independent
routes ($\Omega_\Lambda = 0.6899$ from PSC epoch count; $\Omega_\Lambda = 3\pi/14 =
0.6732$ from three-tape holographic mode count):
(1)~\textbf{PSC epoch-count route}: $\Omega_\Lambda = (\ln 2/N_{\rm gen}\,\pi)\log_2(2000/3) = 0.6899$
from the PSC reflexive-closure orbit count, which admits a six-step first-principles
derivation (CatAD): $|\mathcal{O}_{\rm GTE}| = D^2 N_{\rm fam}^{N_{\rm gen}}/N_{\rm gen}
= (N_{\rm gen}+1)^2\cdot N_{\rm fam}^{N_{\rm gen}}/N_{\rm gen} = 4^2\cdot 5^3/3 = 2000/3$,
with $D=4$ from DPP ($\mathbf{A}_{\rm Lean}$), $N_{\rm fam}=5$ from P01 ($\mathbf{A}_{\rm Lean}$),
and $D^2=16$ as the gravitational orbit discriminant;
(2)~\textbf{Three-tape holographic route}: $\Omega_\Lambda = 3\pi/14 = 0.6732$
from the three-tape mode count $3L$ vs.\ volume $L^3$ (holographic mode
density $3L/L^3 \to 0$ as $L \to \infty$; the vacuum mode fraction gives
$\Omega_\Lambda = \tau \cdot \pi/2 = (3/7)(\pi/2) = 3\pi/14$).
Both routes bracket the observed Planck~2018 value $0.6889 \pm 0.0056$ from above
and below, with zero free parameters ($\mathbf{A/D}$, CatAD).

\textbf{PSC-Adjudication Theorem} (CatAD, P47).  Given the PSC axioms (Two-Layer PSC Theorem,
$\mathbf{A}_{\rm Lean}$), MFRR [T], \lean{op9\_catal\_unconditional} ($\mathbf{A}_{\rm Lean}$),
DPP ($\mathbf{A}_{\rm Lean}$), and gauge spectrum ($\mathbf{A}_{\rm Lean}$), the
cosmological-constant fraction satisfies
$\Omega_\Lambda \in [3\pi/14,\; 0.6899]$ with no free parameters.
The Planck~2018 value $0.6889\pm0.0056$ lies within this bracket ($+0.18\sigma$
from the PSC route).  The EW scale $v_{\rm PSC}$ does not enter $\Omega_\Lambda$
numerically; it enters only the energy density $\rho_\Lambda$ as a kink mass
reference.  The SRRG role is \emph{logical}: $g^* = 1/\varphi$ certifies the
MDL-optimal and PSC-adjudication attractor --- the same fixed point that
selects the EW scale also certifies the CC bracket.  This is strictly
stronger than anthropic fine-tuning: the constraint applies to any PSC-consistent
self-referential physical system, not only observer-containing universes.
The bracket structure is itself a theorem of the adjudication architecture: the
irreducible residual of the PMDL action is the realized diagonal record ledger,
a left-computably-enumerable real of Turing degree exactly $0'$, so
$\Omega_\Lambda$ is limit-computable but not computable, and the two structural
routes are precisely its computable bracket evaluations --- the census-capacity
ceiling above and the carrier-price floor below (CatAD, conditional on a named
record-ledger premise bundle;
P47, P51~\cite{SpivackGTECosmology,SpivackPolynomialTransputation}).

$N_{\rm gen}=3$ is the \emph{unique} integer for which both structural
$\Omega_\Lambda$ routes fall within $5\sigma$ of the Planck~2018 observed value
(all other integers: at least one route deviates by $>35\sigma$; CatAD).
Since PSC selects $N_{\rm gen}=3$ independently through the Two-Layer PSC Theorem,
this constitutes a non-trivial structural consistency of the framework.
A measurement-free counterpart sharpens this: under the derived orientation of
the two routes --- the holographic evaluation prices the MDL-minimal
self-description carrier (the floor), while the census evaluation prices that
carrier's information capacity (the ceiling) --- the admissible dark-energy
interval is non-empty only for $N_{\rm gen} \leq 3$; for every $N \geq 4$ the
floor exceeds the ceiling and no admissible dark-energy fraction exists at all,
with zero observational input.  Combined with the Layer~I bound
$N_{\rm gen} \geq 3$ of the Two-Layer PSC Theorem, this determines
$N_{\rm gen} = 3$ without invoking Presentation Invariance (CatAD, conditional on
the same structural premise bundle;
P47, P48~\cite{SpivackGTECosmology,SpivackGTECompleteFramework}).

The classical cosmological constant vanishes exactly ($\langle V \rangle = 0$ in
the PSC vacuum, Lean-certified), quantum vacuum degeneracy is protected by
$\mathbb{Z}_7$ anomaly-freedom in the pure-$\varphi$ sector
($\mathbf{A}_{\rm Lean}$; see \S\ref{subsec:overview-gravity}),
and one-loop corrections are suppressed by the holographic mode count ($\sim 10^{-42}$).
The same coupling structure resolves the cosmological domain-wall problem: the
$\mathbb{Z}_7$ vacuum symmetry, restored at $T_G \approx 0.70$~GeV after reheating,
produces a transient Kibble wall foam that annihilates via the canonical coupling
bias within ${\sim}10^{-23}$~s --- roughly $700\times$ before BBN --- leaving zero
surviving defect network, with relic signals $\Omega_{\rm GW}h^2 \lesssim
3\times10^{-47}$ at $f \approx 0.02$~Hz and $\Delta N_{\rm eff} \approx 0$
(a parameter-free falsifiable prediction).  Any confirmed cosmological
$\mathbb{Z}_7$ wall relic would falsify the annihilation mechanism and the
associated thermal-history epoch, though not the substrate identification or the
particle-spectrum derivations~\cite{SpivackGTECosmology}.
The paper also derives the CMB spectral index $n_s =
1 - \ln 2/(2\pi^2) = 0.96488$ certified by fourteen zero-sorry Lean theorems,
the CKM CP phase $\delta_{CP} = 68.51^\circ$ ($0.017\%$ from PDG) from an $S_3$
subgroup chain, and the neutrino mass sum $\Sigma m_\nu = 59.4$~meV, all from
zero free parameters with Lean~4 machine-verified steps.
Together, P44--P47 close the programme's open questions on 3+1D emergence, quantum
gravity perturbativity, the Single-Source Principle, and first-principles cosmological
predictions.

Paper~P48~\cite{SpivackGTECompleteFramework} is the synthesis capstone: a 340+ page
monograph presenting the complete PSC-to-cosmology derivation chain in fourteen chapters,
integrating every result from P01--P47 with each step labeled by certification level.
It traces substrate selection ($\mathbb{Z}_7 \times \mathbb{Z}_3$ MDL-minimal),
the $\Phi_{\rm MDL}$ field and particle spectrum, all four SM interactions,
three-tape spacetime and quantum gravity, Born-rule quantum mechanics, and
zero-parameter cosmology, with explicit comparison to eight alternative frameworks
and a complete falsification table.  A dedicated chapter, \emph{The Mathematical
Substrate}, develops the UGP/GTE arithmetic in full --- including the UWCA, the
substrate's native computer --- and gives the normative statement of the
ontological ladder introduced in \S\ref{subsec:overview-substrate}; the Selection
chapter presents the two MDL selections as sequential, the rule selection
consuming the seed selection's orbit
(Direct-Interpolation Lift Theorem, $\CatAL$; full treatment in P49).

Paper~P49~\cite{SpivackGTEPolynomialWolfram} presents the first systematic study of
$p(L,C,R) = C+R-CR-LCR$ over $\mathbb{Z}_7$ as a standalone dynamical system,
provides the Wolfram framing of GTE's MDL selection result (the MDL uniqueness
theorem T96-02 is the selection principle the Wolfram Physics Project has lacked
--- a chain machine-certified link by link, $\CatAL$: the cross-family coding
inequality forcing $k=7$, plus within-class uniqueness via the
Direct-Interpolation Lift Theorem, under which the UGP generation orbit's
total-parity shadow with vacuum transparency forces $p$ uniquely among all $7^8$
multilinear rules, Rule~110 arising as the binary-restriction corollary),
and proves five new algebraic theorems: Invariant Subset Uniqueness ($\CatAL$),
QNR Binary Floor ($\CatAL$), Polynomial Uniqueness ($\CatA$), $f_{\rm MDL}$
non-polynomial ($\CatA$, Schwartz--Zippel), and 40\% DPP cross-tape fraction ($\CatAD$).
It also presents the first $k=7$ polynomial CA spacetime diagrams.

The arithmetic of the polynomial admits a striking unification.  The diagonal
self-consistency equation $p(x,x,x) = x$ factors as $-x(x^2+x-1)$ over
$\mathbb{Z}$, revealing a \textbf{master quadratic} $x^2+x-1$ whose discriminant
equals the family count $N_{\rm fam} = 5$.  This single $\mathbb{Z}$-quadratic has
three arithmetic completions: over $\mathbb{R}$ it gives the SRRG fixed point that
seeds the Higgs VEV ($x = 1/\varphi$, the golden-ratio conjugate;
P27~\cite{SpivackSRRG}); over $\mathrm{GF}(7)$ it is irreducible ($5$ is a
quadratic non-residue mod~$7$), which is exactly why the binary floor $\{0,1\}$ is
unique and Rule~110 is forced; and in polynomial dynamics it controls vacuum
uniqueness ($\mathrm{Fix}(T_n) = \{\text{vacuum}\}$ for every ring size $n$) ---
the electroweak scale and the unique computational floor are the two arithmetic
fibers of one equation (\CatAL;
P49, P48~\cite{SpivackGTEPolynomialWolfram,SpivackGTECompleteFramework}).
Moreover, the two quadratic rings the framework employs --- the golden ring
$\mathbb{Z}[\varphi]$ and the Eisenstein ring $\mathbb{Z}[\omega]$, whose norm-7
residue arithmetic realises the substrate group as
$F_{21} \cong (\mathbb{Z}[\omega]/(3+\omega))^+ \rtimes \mu_3$ --- are the two
quadratic subfields of the single biquadratic compositum
$K = \mathbb{Q}(\sqrt{-3},\sqrt{5})$, of conductor
$15 = N_{\rm gen}\cdot N_{\rm fam}$; membership of the GTE alphabet prime class
reduces to a single Artin-symbol condition on
$K$~\cite{SpivackGTEPolynomialWolfram}.

Paper~P50~\cite{SpivackSpin7Lattice} realises the GTE polynomial as a 7-state
spin lattice model with a phase transition in the 3-state Potts universality
class (falsifiable critical exponents $\nu = 5/6$, $\eta = 4/15$), and
Paper~P51~\cite{SpivackPolynomialTransputation} establishes the polynomial
certificate of transputation: the Three-Level MDL Unification Theorem (theory
selection, field dynamics, and event adjudication as nested instances of one
MDL formula) and the computability classification of the adjudication function
(\S\ref{subsec:overview-qm}).

P50's continuum-limit analysis sharpens the lattice--certificate correspondence
into quantitative lattice field theory.  The thermal ensemble of the spin-7
chain is the maximum-entropy shadow of the CMCA tape, and its spectral gap
obeys an exact chiral gap law $\Delta(\beta) = e^{-3\beta/2}(1 + O(e^{-\beta/2}))$
($\CatAD$, with the exact correction tower): the half-integer slope $3/2$ is the
geometric mean of the directed wall energies of the pair digraph ($\CatAL$) ---
impossible for any parity-symmetric chain, hence a quantitative
parity-violation signature of the substrate chirality, the thermodynamic trace
of the AGL(1,7) chiral reflection --- and the amplitude is exactly 1, derived
by two independent routes.  The gap defines a lattice-QFT calibration
trajectory whose physical point is selected by the Tape Saturation Theorem
($\CatAD$, conditional on the named Compton-Support Criterion, proven minimal
via a Boundary Equivalence Theorem): the matter-sector tape spacing saturates
at $a = \hbar c/\Lambda_{\rm GTE} \approx 0.0972$~fm, with exact fixed-point
signature $\xi(\beta^*) = |\mathbb{Z}_7| = 7$ and lattice dictionary
$aM_{\rm kink} = 1/7$, $am_\varphi = 7/8$ --- the bare mass the P42 lattice
campaigns used, now derived rather than chosen.  The unique empirical
discriminator is the substrate-MC kink correlation length, $\xi_{\rm kink} = 7$
cells (band $6.78$--$7.23 \pm 0.54$); a landing at 7 cells would make the
physical point unconditional.

% ====================================================================
\subsection{Physical Predictions and Cross-Checks}
\label{subsec:overview-predictions}

Table~\ref{tab:overview-predictions} collects key quantitative predictions
against PDG measurements.  Additional physical cross-checks include: neutrino mass
structure and oscillation parameters~\cite{SpivackNeutrino};
the mirror-branch Braid Atlas~\cite{SpivackMirrorBraidAtlas} (P29) derives a parameter-free
dark sector from the $\mathbb{Z}_7$ orbit mirror branch: three dark-lepton generations and
$\mathrm{SU}(3)_{\rm dark}$, with the lightest dark state at $211.9$~MeV, constituting
a full set of falsifiable dark-sector predictions;
nuclear binding
energetics~\cite{SpivackSM}; and P33~\cite{SpivackDeeperConsequences} derives additional arithmetic consequences
of the framework: the EW boson mass staircase, the photon-as-vacuum identification,
prime-lock structure, and the arithmetic closure of SM parameters beyond those already
derived in P01.

\begin{table}[htbp]
\centering
\caption{Selected GTE/UGP predictions vs.\ PDG measurements.
  \textbf{Derived}: analytically derived with zero free parameters.
  $\mathbf{A}_{\rm Lean}$: machine-certified in Lean~4 with zero sorry.}
\label{tab:overview-predictions}
\smallskip
\resizebox{\textwidth}{!}{%
\begin{tabular}{lllll}
\toprule
Quantity & GTE/UGP value & PDG value & Accuracy & Status \\
\midrule
$\sin^2\theta_W$ (tree) & $3/13 = 0.23077$ & $0.23129(4)$ (PDG 2024) & $-13.0\sigma$ & $\mathbf{A}_{\rm Lean}$ \\
$\sin^2\theta_W$ (threshold) & $384729/1664000$ & $0.23122(3)$ (PDG 2022) & $-0.42\sigma$ & $\mathbf{A}_{\rm Lean}$ \\
$\sin^2\theta_W$ (2-loop) & $0.23129154$ & $0.23129(4)$ (PDG 2024) & $+0.038\sigma$ & $\mathbf{B}$ \\
$M_Z$ & $91.914$~GeV & $91.188(2)$~GeV & $+0.80\%$ & $\mathbf{A/D}$ \\
$M_W$ & $80.614$~GeV & $80.369(13)$~GeV & $+0.30\%$ & $\mathbf{A/D}$ \\
$\lambda_{\rm CKM}$ & $9/40 = 0.225$ & $0.22500(67)$ & $0.0\sigma$ & $\mathbf{A}_{\rm Lean}$ \\
$\alpha_s(M_Z)$ & $0.11822$ & $0.1180(9)$ & $+0.24\sigma$ & $\mathbf{A}_{\rm Lean}$ \\
$\alpha_{\rm EM}^{-1}$ & $137$ & $137.036\ldots$ & derived & $\mathbf{A}_{\rm Lean}$ \\
$b_0$ (QCD $\beta_0$) & $7 = |\mathbb{Z}_7|$ & $7$ (exact) & exact & $\mathbf{A}_{\rm Lean}$ \\
$\theta_{\rm QCD}$ & $0$ (exact) & $< 10^{-10}$ & exact & $\mathbf{A}_{\rm Lean}$ \\
$M_{\rm Pl}/m_\tau$ & $F_{21}^{10} \cdot 7^7 / 2$ & PDG ratio & $0.040\%$ & $\mathbf{A/D}$ \\
$\theta_P$ ($\eta$--$\eta'$ mixing) & $-13.08^\circ \pm 3.74^\circ$ & $[-14.3^\circ,\,-10.7^\circ]$ & in range & $\mathbf{A/D}$ \\
$N_{\rm gen}$ & $3$ (derived) & $3$ & exact & $\mathbf{A}_{\rm Lean}$ \\
$\Lambda_{\rm classical}$ & $0$ (exact) & $\approx 10^{-52}$~m$^{-2}$ (quantum) & classical exact & $\mathbf{A/D}$ \\
$n_s$ (CMB tilt) & $1 - \ln(2)/(2\pi^2) = 0.96488$ & $0.9649 \pm 0.0042$ (Planck~2018) & $-0.005\sigma$ & $\mathbf{A/D}$ (cond.) \\
$r$ (tensor-to-scalar) & $0$ (exact) & unmeasured ($\sigma_r \approx 0.001$ at LiteBIRD) & primary falsifiable & $\mathbf{A}$ \\
$y_\tau$ (tau Yukawa) & $1/98 = 1/(2{\times}7^2)$ & SM value & exact structural & CatA \\
$m_\tau$ (from $v_H$, no mass input) & $1776.57$~MeV & $1776.86$~MeV & $0.016\%$ & CatA \\
$m_\mu$ (Koide from $y_\tau$) & $105.6$~MeV & $105.658$~MeV & $0.022\%$ & CatA \\
$m_e$ (Koide from $y_\tau$) & $0.511$~MeV & $0.51100$~MeV & $0.023\%$ & CatA \\
$g_{hKK}$ (kink--Higgs coupling) & $4/7^4 = M_{\rm kink}/(v_H/\sqrt{2})$ & --- & structural & CatA \\
$\Omega_\Lambda$ (PSC route, P47) & $(\ln 2/3\pi)\log_2(2000/3) = 0.6899$ & $0.6889 \pm 0.0056$ (Planck~2018) & $+0.18\sigma$ & CatAD \\
$\Omega_\Lambda$ (holographic, P47) & $3\pi/14 = 0.6732$ & $0.6889$ (Planck~2018) & bracket & CatAD \\
$\Omega_\Lambda$ PSC bracket & $[3\pi/14,\; 0.6899]$; $N_{\rm gen}{=}3$ unique & $0.6889$ (Planck~2018) & within bracket & CatAD \\
$C_{\rm Gorard}$ & $N_{\rm spatial}/(2D^2) = 3/32$ exact & --- & exact & $\mathbf{A}_{\rm Lean}$ (\lean{c\_gorard\_eq\_n\_spatial\_over\_two\_dsq}) \\
$\delta_{CP}$ (CKM $S_3$ chain, P47) & $\pi/2 - 3/8 = 68.51^\circ$ & $68.517^\circ$ (PDG) & $0.017\%$ & CatAD \\
$\Sigma m_\nu$ (P47) & $59.4$~meV & $< 120$~meV (Planck) & within bound & CatAD \\
\midrule
\multicolumn{5}{l}{\textit{New machine-certified results (2026)}} \\
\midrule
PMNS $\sin^2\theta_{12}$ (P21) & $4/13 = 0.3077$ & $0.308\pm0.011$ (NuFIT~6.0 IC24 NH) & $+0.01\sigma$ & $\mathbf{A}_{\rm Lean}$ \\
PMNS $\sin^2\theta_{23}$ (P21) & $19/42 = 0.4524$ & $0.470\pm0.013$ (NuFIT~6.0 IC24 NH; $43.3^\circ$) & $-1.35\sigma$; within $3\sigma$ & $\mathbf{A}_{\rm Lean}$ \\
PMNS $\sin\theta_{13}$ (P21) & $11/73 = 0.1507$ & $0.1489\pm0.0007$ (NuFIT~6.0 IC24 NH) & $+1.0\sigma$ & $\mathbf{A}_{\rm Lean}$ \\
Jarlskog $J_{\rm GTE}$ (P21) & $-1.468{\times}10^{-2}$ (from orbit ratios) & $1.78{\times}10^{-2}$ (NuFIT~6.0 IC24 NH) & see P21 & $\mathbf{A}_{\rm Lean}$ \\
Higgs mass $m_H$ (P27, P35) & $125.2499$~GeV (SRRG-corrected $\lambda$) & $125.20\pm0.11$~GeV & $-0.0003\sigma$ & $\mathbf{A}_{\rm Lean}$ \\
$\eta_B$ (P47, P42; FKTT) & $6.109\times10^{-10}$ & $(6.10\pm0.06)\times10^{-10}$ & $+0.15\sigma$ & $\mathbf{A}_{\rm Lean}$ \\
$H_0$ (P47) & $67.95$~km/s/Mpc (first-principles via GTB kink-overlap cert; epoch-dependent) & $67.66\pm0.42$~km/s/Mpc & $+0.7\sigma$ & CatAD \\
Quantum kink mass $M^Q$ (P38, P42) & $281 \pm 21$~MeV (interface dim-reg, two routes) & --- & see P42 & CatA \\
FGCI (P01) & $b_{L_1}=73$ doubly determined; unique at $N_c=3$ & --- & exact & $\mathbf{A}_{\rm Lean}$ \\
$\Omega_{\rm DM} h^2$ (P29) & $0.11994$ (via $D_{\rm top} = e^{-1/N_c}$, machine-certified \CatAL) & $0.1200\pm0.0012$ & $-0.044\%$ & CatAD \\
\bottomrule
\end{tabular}
}% end resizebox
\end{table}

% ====================================================================
\subsection{Experimental Programme}
\label{subsec:overview-experimental}

The framework makes six sharp falsifiable predictions testable by currently
running or near-term precision experiments, collected in
Table~\ref{tab:overview-experimental}.  Proton stability provides the sharpest
topological test: the Lean~4 theorem \lean{proton\_decay\_dim4\_forbidden}
($\mathbf{A}_{\rm Lean}$, zero sorry) proves that all
\emph{dimension-4} baryon-number-violating vertices are forbidden by
$\mathbb{Z}_7$ winding conservation --- a direct test of the interaction
skeleton theorem.  Any confirmed dimension-4 proton decay would falsify the
$\mathbb{Z}_7$ winding conservation law; the programme's other derivations
(mass spectrum, cosmological predictions, three-generation proof) would be
unaffected.  Notably, standard GUT dimension-6 proton decay
(e.g.\ $p \to e^+\pi^0$, the primary mode searched at Hyper-Kamiokande and DUNE)
is \emph{not} forbidden by the framework and, if observed, would be
completely consistent with the programme's predictions.

\begin{table}[htbp]
\centering
\caption{Six falsifiable predictions of the GTE/UGP programme with active or
  near-term experimental tests.  The ``Falsifying Condition'' column specifies
  the observation that, if confirmed at the stated significance, would require
  structural revision of the framework.}
\label{tab:overview-experimental}
\smallskip
\resizebox{\textwidth}{!}{%
\begin{tabular}{@{}p{3.0cm}p{5.2cm}p{3.4cm}p{1.8cm}p{4.6cm}@{}}
\toprule
Observable & GTE Prediction & Experiment(s) & Timeline & Falsifying Condition \\
\midrule
Proton stability (dim-4) &
  No \emph{dimension-4} baryon-number-violating decay; $\mathbb{Z}_7$ winding
  conservation forbids all dim-4 SU(5) vertices
  (\lean{proton\_decay\_dim4\_forbidden}, $\mathbf{A}_{\rm Lean}$; P22, P33).
  \emph{Note:} standard GUT dim-6 modes ($p\!\to\!e^+\pi^0$,
  $p\!\to\!\bar\nu K^+$) are \textbf{not} forbidden and would be consistent
  with this framework if observed. &
  Hyper-Kamiokande, DUNE &
  2027$^+$ &
  Any confirmed \emph{dimension-4} baryon-violating event (not dim-6 GUT
  modes such as $p\!\to\!e^+\pi^0$, which are not forbidden) \\[3pt]
GTE-P7 dark particle &
  Neutral pseudoscalar at $211.9$~MeV, $Q = 0$;
  mirror-branch $\mathbb{Z}_7$ prediction, machine-certified in Lean~4 (P02, P29) &
  Belle~II (mono-photon / missing energy) &
  Running &
  Non-observation at $50\,\mathrm{ab}^{-1}$ \\[3pt]
PMNS Dirac CP phase &
  $\delta_{CP}^{\rm PMNS} = 205.71^\circ = \tfrac{4}{7}\times360^\circ$;
  from $\mathbb{Z}_7$ winding arithmetic (P45) &
  DUNE, Hyper-Kamiokande &
  2027$^+$ &
  Measurement outside $205^\circ \pm 15^\circ$ \\[3pt]
Tensor-to-scalar ratio &
  $r = 0$ exactly; no primordial gravitational waves in the
  GTE inflation-free cosmology (P44, P45; $\mathbf{A}_{\rm Lean}$) &
  LiteBIRD &
  ${\sim}2028$ &
  $r > 0.01$ at $> 3\sigma$ \\[3pt]
Neutrino mass sum &
  $\Sigma m_\nu = 59.4$~meV; normal mass ordering;
  individual spectrum $\{0.68,\, 8.6,\, 50.1\}$~meV (P21, P47; analytically derived) &
  EUCLID, CMB-S4 &
  2026$^+$ &
  $\Sigma m_\nu < 40$~meV or $> 80$~meV \\[3pt]
Dark-energy EOS &
  $w = -1$ exactly; cosmological-constant behavior forced by
  PSC boundary condition, Weinberg no-go evaded (P47; analytically derived;
  note: DESI~2024 data favour $w \approx -0.95$) &
  DESI, EUCLID &
  Ongoing &
  $w \neq -1$ at $> 3\sigma$ \\
\bottomrule
\end{tabular}
}% end resizebox
\end{table}

% ====================================================================
\subsection{Open Questions}
\label{subsec:overview-open}

The programme is mature but not complete.  The following structural gaps remain open
as of this writing.  A number of questions previously listed as open have been closed
by P44--P47 and subsequent analytical work: 3+1D emergence has been derived (DPP, $\CatAL$);
the quantum gravity sector is perturbatively complete (P44); the Single-Source Principle
is established ($\CatAL$); $r = 0$ ($\mathbf{A}$) and $\delta_{CP} = 68.51^\circ$ (CatAD, P47) are derived as
falsifiable predictions; the Newtonian gravity force law $F\propto r^{-2}$ has been
established at CatAD (three equivalent derivations); the Born rule is
independently confirmed at CatAD via the Page--Wootters $\tau_c$ mechanism;
the $\Omega_\Lambda$ bracket $[3\pi/14,\;0.6899]$ is CatAD with no free parameters (PSC-Adjudication Theorem, P47);
the SRRG--MDL equivalence OP9 is CatAL at the fixed point and CatAD for the full biconditional;
the $\mathbb{Z}_7$ anomaly-freedom is $\mathbf{A}_{\rm Lean}$;
the \textbf{unconditional MDL uniqueness of $\mathbb{Z}_7\times\mathbb{Z}_3$}
is $\mathbf{A}_{\rm Lean}$ (\lean{mdl\_ca\_rule\_coding\_closed},
\lean{z5\_fmdl\_no\_psc\_kink\_orbits}, \lean{mdl\_total\_z7z3\_strictly\_beats\_z5z3},
all zero sorry; the GTE polynomial over $\mathrm{GF}(5)$ has zero PSC kink orbits,
giving total MDL gap $3+0 < 5+6$ with no SM input at any step);
and cosmological predictions for $\Omega_\Lambda$, $n_s$, and the CKM CP phase have
been derived from first principles in P47.

\begin{itemize}[leftmargin=1.5em,itemsep=0.25em]
  \item \textbf{Explicit glider-to-fermion map.}
    The existence of a canonical map from specific Rule~110 glider patterns to
    specific SM fermions is confirmed via the Algebraic Lifting Theorem
    ($\mathbf{A}_{\rm Lean}$).  The mesoscopic orbital form is machine-certified:
    the $f_{\rm MDL}$ forward chain $\mathrm{gen}_1 \to \mathrm{gen}_2 \to
    \mathrm{gen}_3 \to \mathrm{vacuum}$ assigns a unique orbital position to each
    SM fermion generation (\lean{zgm\_mes\_orbital\_discriminant},
    \lean{ZGMMesInvariant.lean}, $\mathbf{A}_{\rm Lean}$).  The hard form ---
    explicit one-to-one identification of each spatial Rule~110 glider species
    with a specific fermion --- remains the deepest open particle-ontology question.

  \item \textbf{Gluon first-principles discrimination (colored $W_B{=}0$ sector).}
    The colorless $W_B{=}0$ sector is fully discriminated: photon, neutrino, $Z$,
    and Higgs are pairwise distinguished by a three-level first-principles labeling
    (\lean{z7\_zero\_sector\_unified\_discriminant},
    \lean{Z7ZeroSectorDiscriminant.lean}, $\mathbf{A}_{\rm Lean}$, zero sorry).
    Gluons also carry $W_B = 0$ but are colored ($Q_\chi \neq 0$) and are excluded
    from the colorless hierarchy; a complete CA-dynamics account of all eight gluon
    species from the $F_{21}$ orbit structure remains open.

  \item \textbf{Newton's constant in SI units.}
    The ratio $G_N = m_\tau^2/M_{\rm Pl}^2$ is derived to 0.040\% ($\mathbf{A/D}$).
    A purely first-principles derivation of $G_N$ in SI units from the $\tau_c$
    dimensional mechanism --- without using $m_\tau$ as an input --- remains open.

  \item \textbf{Quantum cosmological constant.}
    Significant progress has been made.
    The classical $\Lambda = 0$ is Lean-certified ($\mathbf{A/D}$).
    Quantum vacuum degeneracy across all 7 $\mathbb{Z}_7$ sectors is an
    $\mathbf{A}_{\rm Lean}$ theorem (\lean{z7\_vacuum\_sectors\_equiprobable};
    pure-$\varphi$ sector, $\mathbb{Z}_7$-periodic observables).
    The PSC-adjudication bracket $\Omega_\Lambda \in [3\pi/14,\;0.6899]$ is
    established at CatAD with no free parameters (P47), and one-loop corrections
    are suppressed by $\sim 10^{-42}$ via holographic non-renormalization (CatAD).
    \textbf{All computable quantum corrections to $\Lambda$ vanish identically in
    the pure-$\varphi$ sector} --- machine-certified by
    \lean{lambda\_all\_order\_zero\_pure\_phi} (P48, $\mathbf{A}_{\rm Lean}$,
    zero sorry).  This closes OP3 at $\mathbf{A}_{\rm Lean}$ for the pure-field
    sector.  What remains open is the full-Lagrangian sector (including
    $V_{\rm coupling}$): a complete non-renormalization proof to all loop orders
    in the interacting theory (CatD).

\end{itemize}

% ====================================================================
\subsection{Summary and Significance}
\label{subsec:overview-summary}

The UGP Physics programme has established, through a combination of Lean~4 machine-checked
proofs and precision numerical verification, that the $\mathbb{Z}_7$-symmetric
Klein--Gordon field $\Phi_{\rm MDL}$ --- uniquely selected as the physical
continuum substrate by MDL minimality and PSC self-consistency from first
principles, formalised as the Möbius triple $(A,e,[D])$ in P34~\cite{SpivackGTEMobius},
and algebraically certified by its discrete counterpart $f_{\rm MDL}$ and the CMCA
(P41/P45) --- generates the quantitative structure of the Standard Model particle spectrum,
all three gauge coupling constants, the electroweak mixing angle, the CKM matrix,
the full QCD sector including colour confinement and asymptotic freedom, Turing
universality, quantum mechanics with its Born rule, exact Lorentz invariance, the
Einstein field equations with Newton's constant, the holographic entropy formula
$S(A) = \mathrm{Area}[\Gamma_{\rm min}]/(4G)$, Hawking temperature unmodified by
$m_\varphi = m_\tau$, the MDL-selected flat cosmological initial state that
dissolves the need for inflation, and the first-principles cosmological predictions
of P47: the dark-energy fraction ($\Omega_\Lambda = 0.6899$ and
$3\pi/14 = 0.6732$, bracketing Planck~2018 within the PSC-adjudication band $[3\pi/14,\;0.6899]$ at CatAD), the CMB spectral index ($n_s = 0.96488$, $-0.005\sigma$),
the CKM CP phase ($\delta_{CP} = 68.51^\circ$, $0.017\%$ from PDG), and the neutrino
mass sum ($\Sigma m_\nu = 59.4$~meV).
Recent machine-certified results additionally establish, all with zero sorry:
the three PMNS mixing angles ($\sin^2\theta_{12} = 4/13$, $\sin^2\theta_{23} = 19/42$,
$\sin\theta_{13} = 11/73$; $\chi^2 = 0.564$, zero free parameters) and the
Jarlskog invariant $J_{\rm GTE} = -(8184\sqrt{437}/5057221)\sin(\pi/7)$;
the corrected Higgs mass $m_H = 125.2499$~GeV ($-0.0003\sigma$ PDG~2022) via the
SRRG identity $2c_H+1=27=N_{\rm gen}^3$ (\CatAL);
the baryon asymmetry $\eta_B = 6.109\times10^{-10}$ ($+0.15\sigma$) via the FKTT mechanism
(\CatAL, zero sorry, zero axioms; \lean{kink\_top\_coupling\_eq\_eps\_FN}); and
the Frobenius-Generation Coincidence Identity confirming that
$b_{L_1} = 73$ is doubly determined, witnessing consistency of the algebraic and
dynamical structures at $N_c = 3$ (\CatAL; \lean{FrobeniusChain.lean}).
No free parameters are introduced at any stage; the only true external input is one
energy scale ($G_F$, which fixes $v_H = (\sqrt{2}\,G_F)^{-1/2}$).  The tau-lepton
Yukawa $y_\tau = 1/98$ and the structural Koide relation ($\theta = 2/9$ from $N_c = 3$)
then derive all three charged-lepton masses from $v_H$ without any fermion mass anchor.

The ugp-lean library currently comprises over 300 Lean~4 modules, including zero-sorry coverage of
\lean{FrobeniusChain.lean} (17 theorems; FGCI), \lean{MassRelations/NeutrinoSector.lean}
(PMNS angles, Jarlskog, FN rank-1 theorem), \lean{MassRelations/HiggsQuartic.lean}
($m_H = 125.2499$~GeV, $2c_H+1=N_{\rm gen}^3$),
\lean{UCLMassOrdering.lean} (fermion mass ordering),
and \lean{Gravity/\allowbreak{}YukawaOverlap\allowbreak{}Exponent.lean} ($\eta_B$ overlap exponent).
Every Category-A physics theorem in the library has axiom closure exactly
$\{\texttt{propext},\, \texttt{Classical.choice},\, \texttt{Quot.sound}\}$ ---
the standard Mathlib signature.  No UGP-specific axioms appear in any
Category-A physics theorem.

The programme's distinguishing feature is not merely quantitative agreement but
logical necessity: the constants are not fitted but \emph{forced}.  The number three
is forced for the generation count; the fraction $3/13$ is forced for the Weinberg
angle; the ratio $7$ is forced for the QCD $\beta$-function coefficient; the
exponential hierarchy between $M_{\rm Pl}$ and $m_\tau$ is forced by $F_{21}$
orbit combinatorics; $r = 0$ and $\delta_{CP} = 68.51^\circ$ follow from the three-tape clock architecture and CKM orbit arithmetic (P47);
and the dark-energy fraction, CMB tilt, CKM CP phase, and neutrino masses all follow
from the PSC holographic structure of P47.
If a future precision measurement displaces any of these
predictions, the framework's foundational principle --- PSC (Perfect Self-Containment), with MDL
as its operational expression, and $F_{21}$ symmetry as the unique admissible
internal structure --- provides unambiguous targets for revision.

The 52 papers of the UGP Physics programme (P01--P52), the ugp-lean Lean~4 library, and the
computational infrastructure are all archived on Zenodo and catalogued at
\url{https://www.novaspivack.com/research}~\cite{SpivackResearchProgrammes}.

\medskip
\noindent\textbf{One field.  One foundational principle --- self-containment --- from
which all of physics follows.}

